% concatenated from On_Generalizations_of_MM_and_PSap.tex
\documentclass[12pt]{article}

\usepackage[a4paper,includeheadfoot,margin=2.54cm]{geometry}
\usepackage[english]{babel}
\usepackage{amsthm}
\usepackage{amsmath}
\usepackage{amssymb}
\usepackage[hyphens]{url}

\usepackage{multirow,graphicx} 
\usepackage[all]{xy}
\usepackage{xcolor}

\theoremstyle{plain}
\newtheorem{theorem}{Theorem}
\newtheorem{lemma}{Lemma}
\newtheorem{proposition}[lemma]{Proposition}
\newtheorem{corollary}[lemma]{Corollary}
\newtheorem{conjecture}{Conjecture}
\newtheorem{fact}[lemma]{Fact}
\newtheorem{example}{Example}[section]
\theoremstyle{definition}
\newtheorem{definition}{Definition}
\newtheorem{question}{Question}
\newtheorem{observation}{Observation}
\newtheorem{problem}{Problem}
\newtheorem{tab}{Table}
\newtheorem{remark}{Remark}
\newtheorem{result}{Result}

\newcommand{\fn}{\mathbb{F}_{2^n}}
\newcommand{\f}{\mathbb{F}_{2}}
\newcommand{\F}{\mathbb{F}}
\newcommand{\Z}{\mathbb{Z}}
\newcommand{\V}{\mathbb{V}}
\newcommand{\Q}{\mathbb{Q}}
\renewcommand{\to}{\rightarrow}
\newcommand{\Tnm}{{\rm Tr^n_m}}
\newcommand{\Tr}{{\rm Tr}}

\newcommand{\new}[1]{{\color{red} #1}}

\begin{document}

\title{On Generalizations of Maiorana--McFarland and $\mathcal{PS}_{ap}$ Functions}

\author{ 
	Sezel Alkan$^{1}$, Nurdag\"{u}l Anbar$^{1}$, Athina Avrantini$^{2}$,\\Erroxe Etxabarri-Alberdi$^{3}$, Tekg\"{u}l Kalayc\i$^{4}$, Beatrice Toesca$^{5}$
	\vspace{0.4cm} \\
	\small $^1$Sabanc{\i} University, MDBF, Orhanl\i, Tuzla, 34956 \. Istanbul, Turkey\\
	\small $^2$Department of Mathematics, University of Pennsylvania, Philadelphia, PA 19104-6395, USA
	\\
	\small $^3$Basic Sciences Department,  Mondragon Unibertsitatea, 20500 Arrasate, Gipuzkoa, Spain
	\\
	\small $^4$Institut f\"ur Mathematik, Alpen-Adria-Universit\"at Klagenfurt, Austria \\
	\small $^5$Institute of Mathematics, University of Zurich, Switzerland
	\\
	\small Email: {\tt sezel.alkan@sabanciuniv.edu } \\
	\small Email: {\tt nurdagulanbar2@gmail.com}\\
	\small Email: {\tt athina@sas.upenn.edu }\\
	\small Email: {\tt eetxabarri@mondragon.edu  } \\
	\small Email: {\tt tekgulkalayci1@gmail.com}\\
	\small Email: {\tt beatrice.toesca@math.uzh.ch  }
}
\date{}

	\maketitle
	
	\begin{abstract}
		We study generalizations of two classical primary constructions of Boolean bent functions, namely the Maiorana--McFarland ($MM$) class and the (Desarguesian) partial spread ($\mathcal{PS}_{ap}$) class. 
		
		The construction of bent functions lying outside the completed $MM$ class has attracted considerable attention in recent years. 
		In this direction, we construct families of generalized Maiorana--McFarland bent functions that are not equivalent to any function in the classical $MM$ or $\mathcal{PS}_{ap}$ classes, and hence lie outside their completed classes.
		
		As a second contribution, we investigate the decomposition of generalized $\mathcal{PS}_{ap}$ functions. 
		We prove that when the degree is sufficiently small relative to the size of the underlying finite field, such functions do not, in general, admit a decomposition into bent or semibent functions. 
		Consequently, they cannot be obtained from known secondary constructions based on concatenation.
		
		Finally, we present a secondary construction of Boolean bent functions arising from the concatenation of components of vectorial generalized $\mathcal{PS}_{ap}$ functions.
		
		Our constructions and proofs rely on classical results concerning second-order derivatives of bent functions and their duals. In addition, we employ methods from the theory of algebraic curves and their function fields.
	\end{abstract}
	
	\noindent\textbf{Keywords:} Boolean functions, Bent functions, Concatenation, Decomposition, Maiorana--McFarland functions, Partial spread functions, Algebraic curves, Function fields.
	
	\medskip
	\noindent\textbf{Mathematics Subject Classification (2010):} 11T06, 94A60, 14H05.
	
	\section{Introduction}
	
	Let \(p\) be a prime and let \(\V_n^{(p)}\) be an \(n\)-dimensional vector space
	over the prime field \(\mathbb{F}_p\).
	We fix a non-degenerate inner product
	\(\langle \cdot , \cdot \rangle_n\) on \(\V_n^{(p)}\).
	
	If \(\V_n^{(p)} = \mathbb{F}_p^n\) is the vector space of \(n\)-tuples over
	\(\mathbb{F}_p\), we take the usual dot product \(\langle b, x \rangle_n = b \cdot x\).
	If \(\V_n^{(p)} = \mathbb{F}_{p^n}\) is the finite field of order $p^n$, we define $\langle b, x \rangle_n = \Tr_1^n(bx)$,
	where, for any divisor \(k\) of \(n\), \(\Tr_k^n\) denotes the trace mapping from
	\(\mathbb{F}_{p^n}\) to \(\mathbb{F}_{p^k}\).
	When \(n = 2m\), it is often convenient to identify
	\(\V_n^{(p)} =\mathbb{F}_{p^m} \times \mathbb{F}_{p^m}\), in which case the
	inner product is given by $\langle (u,v), (x,y) \rangle_n = \Tr_1^m(ux + vy)$.
	
	A function \(f : \V_n^{(p)} \to \mathbb{F}_p\) is called a \emph{\(p\)-ary
		function}; in the special case \(p=2\), it is called a \emph{Boolean function}.
	The Walsh transform of \(f\) is the complex-valued function defined by
	\[
	W_f(b)
	=
	\sum_{x \in \V_n^{(p)}}
	\zeta_p^{\,f(x) - \langle b, x \rangle_n},
	\qquad b \in \V_n^{(p)},
	\]
	where \(\zeta_p = e^{2\pi i/p}\) and \(i\) is a primitive fourth root of unity.
	The \emph{extended Walsh spectrum} of \(f\) is the multiset
	\(\{\, |W_f(b)| : b \in \V_n^{(p)} \,\}\).
	
	The function \(f : \V_n^{(p)} \to \mathbb{F}_p\) is called \emph{bent} if $|W_f(b)| = p^{n/2}$ for all $b \in \V_n^{(p)}$.
	Equivalently, \(f\) is bent if and only if, for every nonzero vector
	\(a \in \V_n^{(p)}\), its first-order derivative $D_a f(x) = f(x+a) - f(x)$
	is balanced, i.e., for each \(c \in \mathbb{F}_p\), the equation
	\(D_a f(x) = c\) has exactly \(p^{n-1}\) solutions in \(\V_n^{(p)}\).
	
	In the Boolean case \(p=2\), the Walsh transform of a bent function \(f\)
	satisfies $W_f(b) = 2^{n/2} (-1)^{f^{\ast}(b)}$,
	where \(f^{\ast} : \V_n^{(2)} \to \mathbb{F}_2\) is a Boolean function, called the
	\emph{dual} of \(f\), which is itself bent.
	Since \(W_f\) is integer-valued in this case, Boolean bent functions exist only
	when \(n\) is even.
	
	In contrast, for odd $p$, bent functions 
	$f : \V_n^{(p)} \to \mathbb{F}_p$ exist for both even and odd values of $n$.
	Their Walsh transform satisfies
	\[
	W_f(b) =
	\begin{cases}
	\pm \zeta_p^{\,f^{\ast}(b)} \, p^{n/2},
	& \text{if \(n\) is even, or if \(n\) is odd and \(p \equiv 1 \mod{4}\)}, \\[1ex]
	\pm i\, \zeta_p^{\,f^{\ast}(b)} \, p^{n/2},
	& \text{if \(n\) is odd and \(p \equiv 3 \mod{4}\)},
	\end{cases}
	\]
	where 
	\(f^{\ast} : \V_n^{(p)} \to \mathbb{F}_p\) is again called the \emph{dual} of
	\(f\); see \cite{KSW85}.
	
	A function \(f : \V_n^{(p)} \to \mathbb{F}_p\) is called
	\emph{\(s\)-plateaued} if there exists an integer \(0 \le s \le n\) such that $\lvert W_f(b) \rvert \in \{\, 0,\, p^{(n+s)/2} \,\}$ for all $ b \in \V_n^{(p)}$.
	In particular, \(f\) is bent if and only if \(s = 0\).
	Moreover, in the Boolean case \(p = 2\), since the Walsh transform \(W_f\) takes
	integer values, the integers \(n\) and \(s\) must have the same parity.
	A Boolean function is called \emph{semibent} if \(s = 1\) when \(n\) is odd, or
	if \(s = 2\) when \(n\) is even.
	
	Let \(F : \V_n^{(p)} \to \V_k^{(p)}\) be a function.
	For each \(\alpha \in \V_k^{(p)} \setminus \{0\}\), the associated \emph{component
		function} of \(F\) is defined by $F_\alpha (x) = \langle \alpha , F(x) \rangle_k$.
	The function \(F\) is called \emph{(vectorial) bent} if all its
	component functions \(F_\alpha \), \(\alpha \in \V_k^{(p)} \setminus \{0\}\), are 
	bent functions.
	In this case, the set of component functions of \(F\), together with the zero
	function, forms a \(k\)-dimensional vector space over \(\mathbb{F}_p\)
	of bent functions.
	
	In terms of the Walsh transform, a function
	\(F : \V_n^{(p)} \to \V_k^{(p)}\) is bent if, for every
	\(\alpha \in \V_k^{(p)} \setminus \{0\}\) and every \(b \in \V_n^{(p)}\),
	\[
	W_{F_\alpha}(b)
	=
	\sum_{x \in \V_n^{(p)}}
	\zeta_p^{\,\langle a, F(x) \rangle_k - \langle b, x \rangle_n},
	\qquad
	\zeta_p = e^{2\pi i/p},
	\]
	has absolute value \(p^{n/2}\). Equivalently, a function \(F : \V_n^{(p)} \to \V_k^{(p)}\) is bent if,
	for every nonzero \(a \in \V_n^{(p)}\), the derivative $D_a F(x) = F(x + a) - F(x)$
	is a balanced function from \(\V_n^{(p)}\) to \(\V_k^{(p)}\).
	
	Extended-affine equivalence is a fundamental notion in the theory of
	(vectorial) functions, as it characterizes transformations that
	preserve important cryptographic properties, including the extended Walsh
	spectrum and, in particular, bentness.
	
	Two functions
	\(F, G : \V_n^{(p)} \to \V_k^{(p)}\) are said to be
	\emph{extended-affine equivalent (EA-equivalent) } if
	\[
	G(x) = L_1\!\bigl(F(L_2(x) + a)\bigr) + L_3(x) + b,
	\]
	where \(L_1 : \V_k^{(p)} \to \V_k^{(p)}\) and
	\(L_2 : \V_n^{(p)} \to \V_n^{(p)}\) are linear permutations,
	\(L_3 : \V_n^{(p)} \to \V_k^{(p)}\) is a linear map, and
	\(a \in \V_n^{(p)}\), \(b \in \V_k^{(p)}\).
	In the special case \(k = 1\), EA-equivalence reduces to the following relation
	between functions \(f, g : \V_n^{(p)} \to \mathbb{F}_p\):
	\[
	g(x) = f(L(x) + a) + \langle c, x \rangle_n + b,
	\]
	where \(L\) is a linear permutation on \(\V_n^{(p)}\),
	\(c \in \V_n^{(p)}\), and \(b\in \mathbb{F}_p\).
	
	A class \(\mathcal{C}\) of \(p\)-ary, respectively Boolean, bent functions is
	called \emph{completed} if it is invariant under EA-. In other words, \(\mathcal{C}\) contains a function \(f\) if and only if
	it contains every function EA-equivalent to \(f\).
	The \emph{completion} of \(\mathcal{C}\), denoted by \(\mathcal{C}^{\#}\), is
	the smallest completed class containing \(\mathcal{C}\).
	
	\smallskip
	
	Bent functions can be obtained via two main types of constructions:
	\emph{primary} and \emph{secondary}.
	Primary constructions are direct algebraic methods that produce bent functions
	from scratch, without relying on previously known examples.
	In contrast, secondary constructions generate new bent functions from one or
	more \emph{known} bent (or related) functions by applying transformations or
	combinations that preserve bentness. While many secondary constructions of bent functions are known (see
	\cite[Section~6.1.16]{Carlet2021}), there are only two classical primary
	constructions, namely the Maiorana--McFarland construction~\cite{mm} and the
	partial spread (Dillon) construction~\cite{dillon}.
	
	It is well known that all bent functions in dimensions up to six belong to the
	completed Maiorana--McFarland class. The first examples of bent
	functions outside the \(MM^{\#}\) class were constructed by Dillon~\cite{dillon},
	who gave explicit examples in eight variables. It was later shown in~\cite{ll},
	using computer search, that the \(MM^{\#}\) class constitutes only a very small
	fraction of all bent functions. More precisely, the number of bent functions in
	the \(MM^{\#}\) class is at most \(2^{72}\), whereas there are approximately
	\(2^{106}\) bent functions in dimension eight. 
	This observation initiated an active line of research focused on the
	construction of bent functions lying outside the \(MM^{\#}\) class; see, for instance, the recent papers
	\cite{Letal2025,PasalicEtAl2024,PasalicEtAl2023,pp}, as well as the survey paper
	\cite{Pasalic26} and the references therein.
	
	While exclusion from the completed Desarguesian partial spread
	\(\mathcal{PS}_{ap}^{\#}\) class—a distinguished subclass of the partial spread
	construction—can be detected using invariants such as the algebraic degree or
	the \(2\)-rank (see~\cite{akm23,wfq}), membership in the \(MM^{\#}\) class is
	characterized by a criterion due to Dillon~\cite{dillon}, formulated in terms
	of second-order derivatives.
	
	In this paper, although we present certain results in the more general setting
	of an arbitrary prime \(p\), our primary focus is on the construction and
	classification of (vectorial) Boolean functions.
	
	The paper is organized as follows. In Section~\ref{sec:MP}, we briefly review two classical constructions: 
	Maiorana--McFarland and Desarguesian partial spread bent functions, 
	and their generalizations for arbitrary characteristic $p$.
	In Section~\ref{sec:MM}, we present a construction of Boolean generalized Maiorana--McFarland functions that do not belong to the \(MM^{\#}\) class. In Section~\ref{sec:decom}, we study the decomposition of Boolean generalized $\mathcal{PS}_{ap}$ functions and show that, in general, they do not arise from secondary constructions of bent or semibent functions via concatenation when their degree is sufficiently small relative to the size of the finite field on which they are defined. Finally, in Section~\ref{concat}, we provide a secondary construction of Boolean bent functions motivated by the concatenation of the components of vectorial $\mathcal{PS}_{ap}$ functions.

	
	\section{Generalized Maiorana--McFarland and generalized $\mathcal{PS}_{ap}$ functions} \label{sec:MP}
	
	In this section, we briefly recall two fundamental primary constructions of
	\(p\)-ary (vectorial) bent functions: the
	Maiorana--McFarland class, introduced independently by Maiorana (unpublished)
	and McFarland~\cite{mm}, and the Desarguesian partial spread class
	\(\mathcal{PS}_{ap}\), introduced by Dillon in his Ph.D.\ thesis~\cite{dillon},
	together with their generalizations.
	
	Throughout, we use the bivariate representation $\V_n^{(p)}= \mathbb{F}_{p^m} \times \mathbb{F}_{p^m}$, $n = 2m$,
	and present the corresponding definitions in this setting.
	
	\paragraph{Maiorana--McFarland (\(MM\)) class.}
	A function $f : \mathbb{F}_{p^m} \times \mathbb{F}_{p^m} \to \mathbb{F}_p$
	is said to belong to the Maiorana--McFarland class if it is of the form
	\[
	f(x,y) = \Tr_1^m\!\bigl(x\,\pi(y)\bigr) + g(y),
	\]
	where \(\pi\) is a permutation of \(\mathbb{F}_{p^m}\) and
	\(g : \mathbb{F}_{p^m} \to \mathbb{F}_p\) is an arbitrary function.
	
	In the vectorial case, a function $F : \mathbb{F}_{p^m} \times \mathbb{F}_{p^m} \to \mathbb{F}_{p^m}$
	is called Maiorana--McFarland if it can be written as
	\[
	F(x,y) = x\,\pi(y) + G(y),
	\]
	where \(G : \mathbb{F}_{p^m} \to \mathbb{F}_{p^m}\) is an arbitrary function.
	
	It is well known that \(f\), and respectively \(F\), is bent if and only if
	\(\pi\) is a permutation of $\mathbb{F}_{p^m}$.
	
	The following characterization of Boolean bent functions, stated in a general form for the completed \(MM\) class, was given by Dillon.
	
	\begin{lemma}{\cite[Dillon's criterion]{dillon}}\label{lem:second-order}
		Let \(n = 2m\). A Boolean bent function $f : \V_n^{(2)} \to \mathbb{F}_2$
		belongs to the \(MM^{\#}\) class if and only if there exists an \(m\)-dimensional vector subspace
		\(\mathcal U \subseteq \V_n^{(2)}\) such that, for all \(a,b \in \mathcal U \), the second-order derivatives
		\begin{equation}\label{eq:second-order-derivative}
			D_a D_b f(x)
			= f(x) + f(x+a) + f(x+b) + f(x+a+b)
		\end{equation}
		vanish identically.
	\end{lemma}
	A subspace \(\mathcal{U} \subseteq \V_n^{(2)}\) for which
	\(D_a D_b f(x) = 0\) for all \(a,b \in \mathcal{U}\) and all
	\(x \in \V_n^{(2)}\) is called an \emph{\(\mathcal{M}\)-subspace} of \(f\)
	(see~\cite{pp}). We note that both the number of \(\mathcal{M}\)-subspaces of
	\(f\) and its maximal dimension, called the \emph{linearity index}
	\(\operatorname{ind}(f)\) of \(f\), are invariant under EA-equivalence.
	For a bent function \(f \colon \V_n^{(2)} \to \mathbb{F}_2\), the linearity
	index satisfies \(1 \le \operatorname{ind}(f) \le n/2\);
	see~\cite[Proposition~5.1]{Pasalic26}. Consequently, a bent function \(f\) on
	\(\V_n^{(2)}\) belongs to the \(MM^{\#}\) class if and only if
	\(\operatorname{ind}(f) = n/2\).
	
	\paragraph{Generalized Maiorana--McFarland (\(GMM\)) class.}
	Let \(m\) and \(k\) be integers with \(1 \le k \le m\).
	For each \(z \in \V_k^{(p)}\), let $f_z : \V_m^{(p)} \to \mathbb{F}_p$
	be a \(k\)-plateaued \(p\)-ary function, and define its Walsh support by
	\[
	\operatorname{supp}(W_{f_z})
	= \{\, b \in \V_m^{(p)} : W_{f_z}(b) \neq 0 \,\}.
	\]
	Assume that the Walsh supports are pairwise disjoint, that is,
	\[
	\operatorname{supp}(W_{f_z}) \cap
	\operatorname{supp}(W_{f_y}) = \emptyset
	\quad
	\text{for all } z,y \in \V_k^{(p)},\; z \neq y.
	\]
	Then the function
	\begin{equation}\label{eq:GMM}
		f : \V_m^{(p)} \times \V_k^{(p)} \to \mathbb{F}_p,
		\qquad
		f(x,z) = f_z(x),
	\end{equation}
	is a \(p\)-ary bent function (see~\cite{cmpGMM} for details).
	
	Observe that, for fixed \(y\), the Maiorana--McFarland function
	\(f(x,y) = \Tr_1^m\!\bigl(x\,\pi(y)\bigr) + g(y)\)
	is affine in \(x\), and hence corresponds to an \(m\)-plateaued \(p\)-ary function that is affine on each coset of \(\mathbb{F}_{p^n} \times \{0\}\).
	Note also that when \(k = m\), the construction~\eqref{eq:GMM} reduces to the
	Maiorana--McFarland class.
	For this reason, functions defined by~\eqref{eq:GMM} are called
	\emph{generalized Maiorana--McFarland functions} if they are affine on the cosets of a \(k\)-dimensional subspace.
	
	A commonly used finite field representation of the \(GMM\) construction is the
	following.
	Consider the function $f : \mathbb{F}_{p^n} \times \mathbb{F}_{p^k} \times \mathbb{F}_{p^k}
	\to \mathbb{F}_p$
	defined by
	\begin{align}\label{eq:genMM}
		f(x,y,z) = f^{(z)}(x) + \Tr_1^k(yz),  
	\end{align}
	where, for each \(z \in \mathbb{F}_{p^k}\), the function
	\(f^{(z)} : \mathbb{F}_{p^n} \to \mathbb{F}_p\) is bent.
	For fixed \(z\), the function \(f(\cdot,\cdot,z)\) is a \(k\)-plateaued
	\(p\)-ary function on
	\(\mathbb{F}_{p^n} \times \mathbb{F}_{p^k}\), and \(f\) is affine on the cosets of \(\{0\} \times \{0\} \times \mathbb{F}_{p^k}\).
	Note that the corresponding Walsh
	supports of \(f^{(z)}\) are pairwise disjoint, in accordance with the \(GMM\) construction.
	
	In~\cite[Theorem~3]{aakm25}, the Walsh transform of \(f\) in~\eqref{eq:genMM} and the explicit form of its dual \(f^{\ast}\) are derived. In particular, the dual function \(f^{\ast}\) is again a \(GMM\) function, given by
	\[
	f^{\ast}(x,y,z)
	=
	\bigl(f^{(y)}\bigr)^{\ast}(x)
	- \Tr_1^k(yz),
	\]
	and satisfies \(f^{\ast\ast} = (f^{\ast})^{\ast} = f\).
	
	\begin{remark}
		A related notion of the generalized Maiorana--McFarland class was introduced
		in~\cite[Section~4.3]{Camion92} as follows. Let
		\(0 \le k \le \tfrac{n}{2}-1\), and set \(r = \tfrac{n}{2} - k\) and
		\(s = \tfrac{n}{2} + k\). A Boolean function $f_{\varphi,h} \colon \V_r^{(2)} \times \V_s^{(2)} \to \mathbb{F}_2$
		of the form
		\begin{equation}\label{eq:MM}
			f_{\varphi,h}(x,y) = \langle x , \pi(y) \rangle_r + h(y),
			\qquad x \in \V_r^{(2)},\; y \in \V_s^{(2)},
		\end{equation}
		is said to be a generalized Maiorana--McFarland function, where
		\(\pi \colon \V_s^{(2)} \to \V_r^{(2)}\) and
		\(h \colon \V_s^{(2)} \to \mathbb{F}_2\) are arbitrary functions;
		see also~\cite{Carlet2004}. Observe that, for fixed $x \in \mathbb{F}_{p^n}$ and $z \in \mathbb{F}_{p^k}$, 
		the function defined in \eqref{eq:genMM} is affine in the variable $y$. 
		Consequently, $f$ is equivalent to a function on $\V_{k}^{(2)} \times\V_{n+k}^{(2)}$
		of the form $f(x,y) = \langle x , \pi(y) \rangle_k + h(y)$,
		where $\pi \colon \V_{n+k}^{(2)} \to \V_k^{(2)}$ and 
		$h \colon \V_{n+k}^{(2)} \to \mathbb{F}_2$.
		
		Observe also that, when $k = 0$, the above construction reduces to the classical Maiorana--McFarland class of bent functions, provided that $\pi$ is a permutation, whereas for
		\(k = \tfrac{n}{2}-1\) the completed class defined by~\eqref{eq:MM} contains all Boolean
		functions on \(\V_n^{(2)}\).
		There is an extensive research devoted to the classification of functions
		of the form~\eqref{eq:MM}, in particular to the characterization of those
		functions that are bent; see the survey paper~\cite{Pasalic26} and the
		references therein.
	\end{remark}
	
	\paragraph{Desarguesian partial spread (\(\mathcal{PS}_{ap}\)) class.}
	Let \(n = 2m\).
	A \emph{(complete) spread} of \(\V_n^{(p)}\) is a collection of
	\(m\)-dimensional \(\mathbb{F}_p\)-subspaces
	\(U_0, U_1, \dots, U_{p^m} \subseteq \V_n^{(p)}\)
	such that $U_i \cap U_j = \{0\}$ for all $0 \le i < j \le p^m$.
	Equivalently, every nonzero element of \(\V_n^{(p)}\) belongs to exactly one of
	the sets \(U_j^{\ast} = U_j \setminus \{0\}\).
	
	It is well known that every function
	\( f : \V_n^{(p)} \to \mathbb{F}_p \)
	constructed as follows is a bent function. For every \( c \in \mathbb{F}_p \),
	the function \( f \) maps the elements of exactly \( p^{m-1} \) of the sets
	\( U_j^{\ast} \), \( 1 \le j \le p^{m} \), to the value \( c \),
	and \( f \) is constant on the subspace \( U_0 \).
	
	Similarly, one obtains vectorial partial spread bent functions
	\( F : \V_n^{(p)} \to \V_m^{(p)} \).
	For every \( c \in \V_m^{(p)} \), the function \( F \) maps the elements of
	exactly one of the sets \( U_j^{\ast} \), \( 1 \le j \le p^{m} \), to the value
	\( c \), and \( F \) is constant on the subspace \( U_0 \).
	
	A well-known example of a spread is the \emph{Desarguesian spread}.
	In the bivariate representation
	\(\V_n^{(p)} \cong \mathbb{F}_{p^m} \times \mathbb{F}_{p^m}\),
	it is given by
	\[
	U = \{\, (0,y) : y \in \mathbb{F}_{p^m} \},
	\qquad
	U_s = \{\, (x,sx) : x \in \mathbb{F}_{p^m} \,\},
	\quad s \in \mathbb{F}_{p^m}.
	\]
	Bent functions arising from the Desarguesian spread admit the explicit
	representation
	\begin{equation}\label{eq:PSap}
		f(x,y) = P\!\left(y\,x^{p^{m}-2}\right),
	\end{equation}
	where
	\(P \colon \mathbb{F}_{p^m} \to \mathbb{F}_p\)
	is a balanced function.
	
	Analogously, vectorial bent functions
	\(F \colon \mathbb{F}_{p^{m}} \times \mathbb{F}_{p^{m}} \to \mathbb{F}_{p^{m}}\)
	arising from the Desarguesian spread are represented as
	\[
	F(x,y) = P\bigl(y\,x^{p^{m}-2}\bigr),
	\]
	where
	\(P \colon \mathbb{F}_{p^{m}} \to \mathbb{F}_{p^{m}}\)
	is a permutation.
	
	Functions obtained from the Desarguesian partial spread construction are called 
	\(\mathcal{PS}_{ap}\) functions.
	It has been shown that every (vectorial) bent function arising from a spread of 
	\(\V_n^{(p)}\) has algebraic degree \((p-1)n/2\); 
	see~\cite{dillon} for the case \(p = 2\) and~\cite{nuwi22} for odd primes \(p\).
	
	\paragraph{Generalized \(\mathcal{PS}_{ap}\) class.}
	The construction of bent functions from spreads can be generalized via the
	notion of \emph{(normal) bent partitions} \cite{nuwi22}, defined as follows:\\
	A partition of \(\V_n^{(p)}\) into an \((n/2)\)-dimensional subspace \(U\) and
	sets \(A_1, A_2, \dots, A_K\) is called a
	\emph{normal bent partition of depth \(K\)} if every function
	\(f :\V_n^{(p)} \to \mathbb{F}_p\) satisfying the following conditions is bent:
	\begin{itemize}
		\item[(i)] For each \(c \in \mathbb{F}_p\), exactly \(K/p\) of the sets
		\(A_j\) are contained in the preimage \(f^{-1}(c)=\{x\in \V_n^{(p)} : f(x)=c\}\).
		\item[(ii)] The function \(f\) is constant on the subspace \(U\).
	\end{itemize}
	In particular, if \(U_0, U_1, \dots, U_{p^m}\) is a spread of \(\V_n^{(p)}\),
	then the sets \(U_0, U_1^{\ast}, \dots, U_{p^m}^{\ast}\), where
	\(U_i^{\ast} = U_i \setminus \{0\}\), form a normal bent partition of depth
	\(p^m\).
	
	The first class of bent partitions that are not equivalent to those arising from spreads 
	was introduced in~\cite{nuwi22} for arbitrary prime $p$, and in~\cite{mp} for $p=2$. Since then, several other primary 
	and secondary constructions of bent partitions have been developed; see, 
	for instance, \cite{aakm25,Anbar2025Preprint,akm22a,wfw}, as well as the survey paper 
	\cite{akm25} and the references therein. 
	
	Among these, the construction presented in~\cite{nuwi22} (and  \cite{mp}) admits an explicit 
	representation and can be described as follows.
	
	Let \(m\), \(k\), and \(e\) be integers such that \(k \mid m\),
	\(e \equiv p^{\ell} \mod{(p^{k}-1)}\), and
	\(\gcd(p^{m}-1, e) = 1\).
	Let \(\eta\) denote the multiplicative inverse of \(e\) modulo \((p^{m}-1)\), that
	is, $\eta e \equiv 1 \mod{(p^{m}-1)}$.
	For \(s \in \mathbb{F}_{p^{m}}\), define the sets
	\begin{equation}\label{eq:U}
		U = \{\, (0, y) : y \in \mathbb{F}_{p^{m}} \,\}, \qquad
		U_s = \{\, (x, s x^{e}) : x \in \mathbb{F}_{p^{m}} \,\},
		\qquad
		U_s^{\ast} = U_s \setminus \{(0,0)\}.
	\end{equation}
	Similarly, define
	\begin{equation}\label{eq:Bgamma}
		V = \{\, (x, 0) : x \in \mathbb{F}_{p^{m}} \,\}, \qquad
		V_s = \{\, (s x^{\eta}, x) : x \in \mathbb{F}_{p^{m}} \,\},
		\qquad
		V_s^{\ast} = V_s \setminus \{(0,0)\}.
	\end{equation}
	For \(\gamma \in \mathbb{F}_{p^{k}}\), define
	\begin{align}\label{eq:Agamma}
		A(\gamma)
		=
		\bigcup_{\substack{s \in \mathbb{F}_{p^{m}} \\
				\Tr_{k}^{m}(s) = \gamma}}
		U_s^{\ast},
		\qquad
		B(\gamma)
		=
		\bigcup_{\substack{s \in \mathbb{F}_{p^{m}} \\
				\Tr_{k}^{m}(s) = \gamma}}
		V_s^{\ast}.  
	\end{align}
	Then the collections
	\begin{equation}\label{eq:Gamma}
		\Omega_1 = \{\,U ,\, A(\gamma) : \gamma \in \mathbb{F}_{p^{k}} \,\},
		\qquad
		\Omega_2 =  \{\, V ,\, B(\gamma) : \gamma \in \mathbb{F}_{p^{k}} \,\}, 
	\end{equation}
	are bent partitions of \(\mathbb{F}_{p^{m}} \times \mathbb{F}_{p^{m}}\) of depth
	\(p^{k}\). Moreover, similarly to the case of \(\mathcal{PS}_{ap}\) bent
	functions, the bent functions
	\(f, g : \mathbb{F}_{p^{m}} \times \mathbb{F}_{p^{m}} \to \mathbb{F}_{p}\)
	obtained from \(\Omega_1\) and \(\Omega_2\), respectively, admit explicit
	representations given by
	\[
	f(x,y)
	=
	P\!\left(\Tr_{k}^{m}\!\left(y x^{-e}\right)\right)
	+ c_0 \bigl(1 - x^{p^{m}-1}\bigr),
	\]
	and
	\[
	g(x,y)
	=
	P\!\left(\Tr_{k}^{m}\!\left(x y^{-\eta}\right)\right)
	+ c_0 \bigl(1 - y^{p^{m}-1}\bigr),
	\]
	where \(P : \mathbb{F}_{p^{k}} \to \mathbb{F}_{p}\) is a balanced function and
	\(c_0 \in \mathbb{F}_{p}\) is a constant. Similarly, the vectorial version of the generalized \(\mathcal{PS}_{ap}\)
	functions from $\mathbb{F}_{p^{m}} \times \mathbb{F}_{p^{m}}$ to $\mathbb{F}_{p^{k}}$
	is defined using a permutation
	\(P \colon \mathbb{F}_{p^{k}} \to \mathbb{F}_{p^{k}}\) and $c_0\in \mathbb{F}_{p^{k}}$.
	
	In the case \(k = m\), the partitions \(\Omega_1\) and \(\Omega_2\)
	are equivalent to the Desarguesian spread.
	Accordingly, they are called \emph{generalized Desarguesian spreads}, and the
	bent functions obtained from them are referred to as
	\emph{generalized \(\mathcal{PS}_{ap}\) bent functions}. For details, we refer to \cite{akm25,nuwi22}.
	
	In general, a generalized $\mathcal{PS}_{{ap}}$ bent function does not
	belong to the completed Maiorana--McFarland class. Moreover, there exist
	generalized $\mathcal{PS}_{{ap}}$ bent functions that lie neither in
	the partial spread class nor in the completed Maiorana--McFarland class.
	For further details, see~\cite[Examples~5.1 and~5.2]{akm23}.
	
	We remark that generalized \(\mathcal{PS}_{ap}\) bent functions were originally
	presented without the additional terms
	\(c_0 \bigl(1 - x^{p^{m}-1}\bigr)\) and
	\(c_0 \bigl(1 - y^{p^{m}-1}\bigr)\), respectively, as in the classical
	\(\mathcal{PS}_{ap}\) construction. While any \(\mathcal{PS}_{ap}\) bent
	function is EA-equivalent to one of the form given in
	Equation~\eqref{eq:PSap}, these additional terms are nevertheless necessary in
	order to obtain a complete set of inequivalent generalized
	\(\mathcal{PS}_{ap}\) bent functions; see~\cite[Example~1]{akm25}.
	
	\section{Construction of \(GMM\) functions outside the \(MM^{\#}\) class} \label{sec:MM}
	
	We begin by examining the generalized Maiorana--McFarland construction given in~\cite{cmpGMM}
	over finite fields.
	Specifically, we consider functions
	\(f : \mathbb{F}_{p^n} \times \mathbb{F}_{p^k} \times \mathbb{F}_{p^k}
	\to \mathbb{F}_p\) defined by
	\begin{equation}\label{GMM}
		f(x,y,z) = f^{(z)}(x) + \Tr_1^k(yz),
	\end{equation}
	where, for each \(z \in \mathbb{F}_{p^k}\), the function
	\(f^{(z)} : \mathbb{F}_{p^n} \to \mathbb{F}_p\) is bent.
	
	To characterize functions lying outside the completed Maiorana--McFarland
	class, we make use of Dillon's criterion; see Lemma \ref{lem:second-order}.
	In this context, it is essential to understand the behavior of second-order
	derivatives under EA-equivalence.
	It is shown in~\cite{akkmpp26} that, for a Boolean function
	\(f \colon \V_n^{(2)} \to \mathbb{F}_2\), the property of having constant
	second-order derivatives is invariant under EA-equivalence. For the convenience
	of the reader, we include the proof here.
	
	\begin{lemma}\label{lem:second}
		Let \(f \colon \V_n^{(2)} \to \F_2\) be a Boolean function, and let
		\(a,b \in \V_n^{(2)}\) be linearly independent. Let
		\(L \colon \V_n^{(2)} \to \V_n^{(2)}\) be a linear permutation, and let
		\(c,d \in \V_n^{(2)}\) and \(e \in \F_2\).
		Then \(D_a D_b f(x)\) is identically equal to \(0\) (respectively, identically equal to $1$)
		for all \(x \in \V_n^{(2)}\) if and only if
		\begin{equation}\label{eq:second}
			D_{L^{-1}(a)} D_{L^{-1}(b)} \left(\ f\bigl(L(x)+c\bigr)
			+ \langle d,x\rangle_n + e \ \right) 
		\end{equation}
		is also identically \(0\) (respectively, identically \(1\)) for all
		\(x \in \V_n^{(2)}\).
	\end{lemma}
	
	\begin{proof}
		Since any affine term \(\langle d,x\rangle_n + e\) has vanishing second-order
		derivatives, we may assume without loss of generality that
		\(g(x) = f(L(x)+c)\). By the linearity of \(L\), it follows that
		\begin{align*}
			& D_{L^{-1}(a)} D_{L^{-1}(b)} g(x) \\
			&\quad
			= g\bigl(x + L^{-1}(a) + L^{-1}(b)\bigr)
			+ g\bigl(x + L^{-1}(a)\bigr)
			+ g\bigl(x + L^{-1}(b)\bigr)
			+ g(x) \\
			& \quad
			= f\bigl(L(x + L^{-1}(a) + L^{-1}(b)) + c\bigr) + f\bigl(L(x + L^{-1}(a)) + c\bigr) \\
			&\qquad \qquad \qquad
			+ f\bigl(L(x + L^{-1}(b)) + c\bigr)
			+ f\bigl(L(x) + c\bigr) \\
			& \quad
			= f\bigl(L(x) + a + b + c\bigr)
			+ f\bigl(L(x) + a + c\bigr)
			+ f\bigl(L(x) + b + c\bigr)
			+ f\bigl(L(x) + c\bigr).
		\end{align*}
		Setting \(y = L(x) + c\), we obtain
		\[
		D_{L^{-1}(a)} D_{L^{-1}(b)} g(x)
		= f(y+a+b) + f(y+a) + f(y+b) + f(y)
		= D_a D_b f(y).
		\]
		Since the map \(x \mapsto L(x) + c\) is a permutation of \(\V_n^{(2)}\), the
		assertion follows.
	\end{proof}
	
	\begin{remark}\label{rem:fcircL}
		Let \(f \colon \V_n^{(2)} \to \F_2\) be a Boolean function, and let
		\(L \colon \V_n^{(2)} \to \V_n^{(2)}\) be a linear permutation. Then, by Lemma~\ref{lem:second},
		\(\mathcal{W} \subseteq \V_n^{(2)}\) is an \(\mathcal{M}\)-subspace of \(f\) if and only if
		\(L^{-1}(\mathcal{W})\) is an \(\mathcal{M}\)-subspace of \(f \circ L\).
	\end{remark}
	
	We now analyze the second-order derivative of the function
	\(f : \mathbb{F}_{2^n} \times \mathbb{F}_{2^k} \times \mathbb{F}_{2^k}
	\to \mathbb{F}_2\) defined in Equation~\eqref{GMM}.
	Recall that
	\[
	f(x,y,z) = f^{(z)}(x) + \Tr_1^k(yz).
	\]
	Let
	\(\nu_1 = (u_1,v_1,w_1)\) and
	\(\nu_2 = (u_2,v_2,w_2)\)
	be two distinct nonzero elements of
	\(\mathbb{F}_{2^n} \times \mathbb{F}_{2^k} \times \mathbb{F}_{2^k}\). Then
	\begin{align}\label{eq:secondder}
		&D_{\nu_1}D_{\nu_2}f(x,y,z)
		= f(x,y,z)
		+ f(x+u_1,y+v_1,z+w_1)
		+ f(x+u_2,y+v_2,z+w_2)\\ \nonumber
		&\qquad\qquad 
		+ f(x+u_1+u_2,y+v_1+v_2,z+w_1+w_2)\\[0.5em] \nonumber
		&=f^{(z)}(x)
		+ f^{(z+w_1)}(x+u_1)
		+f^{(z+w_2)}(x+u_2)
		+ f^{(z+w_1+w_2)}(x+u_1+u_2)
		+ \Tr_1^k\!\bigl(w_1v_2+w_2v_1\bigr).
	\end{align}
	
	\begin{remark}\label{rem:Msub}
		In the case \(w_1 = w_2 = 0\), Equation~\eqref{eq:secondder} yields
		\begin{equation}\label{eq:Dfy}
			D_{\nu_1} D_{\nu_2} f(x,y,z)
			=
			f^{(z)}(x)
			+ f^{(z)}(x+u_1)
			+ f^{(z)}(x+u_2)
			+ f^{(z)}(x+u_1+u_2).
		\end{equation}
		In particular, $D_{\nu_1} D_{\nu_2} f(x,y,z)
		=
		D_{u_1} D_{u_2} f^{(z)}(x)$,
		and hence \(D_{\nu_1} D_{\nu_2} f(x,y,z)\) is independent of \(y\).
		Consequently, if the functions \(f^{(z)}\) are \(MM\) functions
		sharing a common \(\mathcal{M}\)-subspace of dimension \(n/2\), then, by
		Dillon's criterion, the function \(f(x,y,z)\) is itself \(MM\). 
				
		Indeed, it suffices to note that \(f(x,y,z)\) admits an \(\mathcal{M}\)-subspace
		of dimension \(k + n/2\).
		Let \(\mathcal{V}\) denote the common \(\mathcal{M}\)-subspace of the functions
		\(f^{(z)}\), with \(\dim \mathcal{V} = n/2\).
		Then, by Equation~\eqref{eq:Dfy}, the set $\mathcal{V}  \times \mathbb{F}_{2^k} \times \{0\}$
		forms an \(\mathcal{M}\)-subspace of \(f\) of dimension \(k + n/2\), which
		establishes the claim.
	\end{remark}
	
	Remark~\ref{rem:Msub} indicates that, in order to construct a \(GMM\) function from \(MM\) functions that does
	not belong to the completed Maiorana--McFarland class, it is necessary to
	combine \(MM\) functions that do not share a common \(\mathcal{M}\)-subspace of
	dimension \(n/2\).
	
	Motivated by this observation, we consider the construction of \(GMM\) functions
	from \(MM\) functions that share only \textit{trivial} \(\mathcal{M}\)-subspaces, i.e.,
	subspaces of dimension at most one.
	Our approach begins with an \(MM\) function on \(\V_n^{(2)}\) admitting a unique
	\(\mathcal{M}\)-subspace \(\mathcal{W}\) of dimension \(n/2\), with the
	additional property that every nontrivial \(\mathcal{M}\)-subspace is contained in
	\(\mathcal{W}\).
	By applying a suitably chosen linear permutation to this function and
	embedding it into a higher-dimensional \(GMM\) framework, we obtain a \(GMM\) function
	that admits only \(\mathcal{M}\)-subspaces of dimension strictly smaller than
	those allowed for \(MM\) functions.
	
	The following proposition formalizes this construction.
	
	\begin{proposition}\label{thm:main}
		Let \(h\) be an \(MM\) function on \(\V_n^{(2)}\) 
		having a unique 
		\(\mathcal{M}\)-subspace \(\mathcal{W}\) of dimension \(n/2\), such that if $\mathcal{U}$ is a nontrivial
		\(\mathcal{M}\)-subspace of $h$, then $\mathcal{U}\subseteq \mathcal{W}$. Let  \(\widetilde{\mathcal{W}}\) be a complementary subspace of \(\V_n^{(2)}\) satisfying 
		\(\V_n^{(2)} =\mathcal{W} \oplus \widetilde{\mathcal{W}}\) (i.e., $\mathcal{W} \cap \widetilde{\mathcal{W}}=\{0\}$ and \(\V_n^{(2)} =\mathcal{W} + \widetilde{\mathcal{W}}\)), and let \(L\) be a linear permutation of \(\V_n^{(2)}\) with 
		\({L}( \widetilde{\mathcal{W}} ) = \mathcal{W}\).
		
		For each \(z\in\mathbb{F}_{2^k}\), define  
		\begin{align}\label{eq:g}
			f^{(z)}(x) =
			\begin{cases}
				h(x), & \text{if } z \in V,\\[6pt]
				(h\circ L)(x), & \text{if } z \not\in V,
			\end{cases}
		\end{align}
		where $	V = \bigl\{\, z \in \mathbb{F}_{2^k} : \Tr_1^k(z)=0 \,\bigr\}$.
		
		With this choice of \(f^{(z)}\) in \eqref{eq:g}, if \(n > 2k+4\), 
		then the function $	f:\V_n^{(2)} \times \mathbb{F}_{2^k} \times \mathbb{F}_{2^k} \to \mathbb{F}_{2}$, defined by $f(x,y,z) = f^{(z)}(x) + \Tr_1^k\!\bigl(y z\bigr)$, is not (equivalent to) an \(MM\) function. 
		Equivalently, if \(\mathcal{U}\) is an \(\mathcal{M}\)-subspace of \(f\), then 
		\(\dim(\mathcal{U}) < n/2 + k\).
	\end{proposition}
	
	
	\begin{proof}
		We recall that, by Remark~\ref{rem:fcircL}, the subspace
		\(L^{-1}(\mathcal{W}) = \widetilde{\mathcal{W}}\) is an \(\mathcal{M}\)-subspace of
		\(h \circ L\) of dimension \(n/2\).
		Moreover, by assumption, \(\widetilde{W}\) is the unique
		\(\mathcal{M}\)-subspace with the property that every nontrivial
		\(\mathcal{M}\)-subspace of \(h\circ L\) is contained in \(\widetilde{W}\).
		Equivalently, for {distinct} nonzero
		\({a},{b}\in \V_n^{(2)}\), the equality $D_{{a}}D_{{b}} \, h (L (x))=0$ holds for all $x\in \V_n^{(2)}$
		if and only if \({a},{b}\in \widetilde{W}\).
		
		\medskip
		
		Let \(\mathcal{U}\) be an \(\mathcal{M}\)-subspace of \(f\).
		We derive an upper bound on the dimension \(\dim(\mathcal{U})\) of $\mathcal{U}$ in two steps.
		
		\medskip
		\noindent\textbf{Step 1.}
		Define
		\begin{equation}\label{eq:UV}
			\mathcal{U}_V
			=
			\bigl\{\,
			(u,v,w)\in \mathcal{U} : \Tr_1^k(w)=0
			\,\bigr\}.
		\end{equation}
		Let
		\(\nu_1=(u_1,v_1,w_1)\) and \(\nu_2=(u_2,v_2,w_2)\)
		be distinct nonzero elements of \(\mathcal{U}_V\).
		Since \(V\) is a subspace, \(w_1+w_2\in V\).
		
		Assume first that
		\(\Tr_1^k(w_1v_2+w_2v_1)=0\).
		Then for any \(z\in V\),
		Equation~\eqref{eq:secondder} yields
		\begin{align*}
			D_{\nu_1}D_{\nu_2}f(x,y,z)
			&= h(x)+h(x+u_1)+h(x+u_2)+h(x+u_1+u_2) \\
			&= D_{u_1}D_{u_2}h(x).
		\end{align*}
		Hence,
		\(D_{\nu_1}D_{\nu_2}f(x,y,z)=0\) for all \(x\in \V_n^{(2)}\)
		if and only if either
		\(u_1,u_2\) are distinct nonzero elements of the
		\(\mathcal{M}\)-subspace of \(h\), or $(u_1,u_2)\in \{(0,u),(u,0),(u,u):u\in \V_n^{(2)}\}$.
		
		Similarly, for \(z\notin V\), $D_{\nu_1}D_{\nu_2}f(x,y,z)
		= D_{u_1}D_{u_2}\, h\circ L(x)$,
		and vanishing occurs if and only if either
		\(u_1,u_2\) belong to the \(\mathcal{M}\)-subspace of
		\(h\circ L\), or  $(u_1,u_2)\in \{(0,u),(u,0),(u,u):u\in \V_n^{(2)}\}$.
		Since \(\mathcal{W}\cap \widetilde{\mathcal{W}}=\{0\}\),
		we must have $(u_1,u_2)\in \{(0,u),(u,0),(u,u):u\in \V_n^{(2)}\}$.
		
		We define $\mathfrak S \subseteq \mathcal U_V$ to be a subset such that, 
		for any $(u_i,v_i,w_i), (u_j,v_j,w_j)\in \mathfrak S$, one has $\Tr_1^k\!\left(w_i v_j + w_j v_i\right)=0$. In other words,
		\[
		\mathfrak{S}
		=
		\left\{
		\nu=(u,v,w)\in \mathcal{U}_V
		\; :\;
		\Tr_1^k\!\left(w v_i + w_i v\right)=0
		\ \text{for all } (u_i,v_i,w_i)\in \mathfrak{S}
		\right\}.\]
		Choose \(\mathfrak{S}\) of maximal dimension.
		If \(\mathfrak{S}\) contains two elements with nonzero first components,
		then these components must be the same.
		Hence, for a fixed nonzero \(u\in \V_n^{(2)} \),
		\(\mathfrak{S}\) is contained in
		\[
		\{0,u\}\times
		\left\{
		(v_j,w_j)\in \mathbb{F}_{2^k}^2
		\; : \;
		\Tr_1^k(w_j)=0
		\ \text{and}\
		\Tr_1^k(w_j v_i + w_i v_j)=0
		\ \text{for all } i
		\right\},
		\]  
		which is an \(\mathcal{M}\)-subspace of \(f\). 
		
		\smallskip
		
		\noindent\textbf{Case (i):}
		If \(\mathfrak{S}=\{(0,0,0)\}\), then
		we claim that \(\dim(\mathcal{U}_V)\leq 2\).
		Suppose, for the sake of contradiction, that \(\mathcal{U}_V\) contains
		three linearly independent vectors
		\(\nu_i=(u_i,v_i,w_i)\), \(i=1,2,3\).
		Since \(\mathfrak{S}=\{(0,0,0)\}\), we have $\Tr_1^k\!\bigl(w_i v_j + w_j v_i\bigr)=1$ for all $i\neq j$.
		By linear independence, the vectors
		\((u_1+u_2,v_1+v_2,w_1+w_2)\) and
		\((u_3,v_3,w_3)\)
		are distinct nonzero elements of \(\mathcal{U}_V\).
		Moreover, we have $\Tr_1^k\!\bigl((w_1+w_2)v_3+w_3(v_1+v_2)\bigr)=0$.
		Hence,
		\((u_1+u_2,v_1+v_2,w_1+w_2)\) and
		\((u_3,v_3,w_3)\)
		both belong to \(\mathfrak{S}\),
		contradicting the assumption that
		\(\mathfrak{S}=\{(0,0,0)\}\).
		
		\smallskip
		\noindent\textbf{Case (ii):}
		Suppose that \(\mathfrak{S}\neq\{(0,0,0)\}\). If
		\(\mathfrak{S} =\mathcal{U}_V\),
		then $\mathrm{dim}(\mathfrak{S}) =  \mathrm{dim}(\mathcal{U}_V)\leq 2k$.
		
		In the case $\mathfrak{S} \not\subseteq \mathcal{U}_V$, choose
		\(\nu_1=(u_1,v_1,w_1)\in \mathcal{U}_V\setminus \mathfrak{S}\).
		By maximality, there exists
		\(\nu_2=(u_2,v_2,w_2)\in \mathfrak{S}\)  with
		\(\Tr_1^k(w_1v_2+w_2v_1)=1\),
		which forces \(u_2=u\) and \(u_1\notin\{0,u\}\).
		Otherwise,
		\[
		h(x)+h(x+u_1)+h(x+u_2)+h\bigl(x+u_1+u_2\bigr)=0,
		\]
		which would imply that, for any \(z\in V\), $D_{\nu_1}D_{\nu_2}f(x,y,z)=1$.
		
		Conversely, for any \(\nu_i=(u_i,v_i,w_i)\in \mathfrak{S}\) with \(u_i=u\),
		we have $\Tr_1^k\!\bigl(w_1v_i+w_iv_1\bigr)=1$.
		Indeed, if this were not the case, then
		\(\{u,u_1\}\subseteq \mathcal{W} \cap \widetilde{\mathcal{W} }\),
		which contradicts our assumption that
		\(\mathcal{W} \cap \widetilde{\mathcal{W} }=\{0\}\).
		
		We claim that any $\mathcal{M}$-subspace of $f$ contains at most two vectors whose last components have trace zero and whose first components are linearly independent.
		This implies that
		\begin{align*}
			\mathcal{U}_V
			\subseteq
			\langle u_1,u_2\rangle
			\times
			\Bigl\{(v,w)\in \mathbb{F}_{2^k}\times \mathbb{F}_{2^k} :
			\Tr_1^k(w)=0
			\Bigr\},
		\end{align*}
		which is a subspace of dimension \(2k+1\).
		Suppose, for the sake of contradiction, that there exist three vectors \(\nu_i=(u_i,v_i,w_i) \in \mathcal{U}_V\), \(i=1,2,3,\) with linearly independent first components satisfying $\Tr_1^k\!\bigl(w_i v_j+w_j v_i\bigr)=1$ for all $i\neq j$.
		Since \(\{u_1,u_2,u_3\}\) is linearly independent, the vectors
		\((u_1+u_2,v_1+v_2,w_1+w_2)\) and
		\((u_3,v_3,w_3)\)
		are distinct nonzero elements of $\mathcal{U}_V$. 
		Moreover, we have $\Tr_1^k\!\bigl((w_1+w_2)v_3+w_3(v_1+v_2)\bigr)=0$.
		By the argument above, this implies that
		\[
		(u_1+u_2,u_3)\in \{(0,u),(u,0),(u,u):u\in \V_n^{(2)}\}.
		\]
		Since \(u_1\neq u_2\) and \(u_3\neq 0\), we must have \(u_1+u_2=u_3\),
		which contradicts the linear independence of \(\{u_1,u_2,u_3\}\).
		
		Thus, any \(\mathcal{M}\)-subspace of \(f\) consisting of 
		\((u,v,w)\) with \(w\in V\) has dimension at most \(2k+1\). In particular, the dimension of $\mathcal{U}_V$ is at most \(2k+1\).
		
		\medskip
		\noindent\textbf{Step 2:}  
		We now consider the set $\bigl\{\, (u,v,w) \in \mathcal{U} : \Tr_1^k(w)=1 \,\bigr\}$. 
		Equivalently, for a fixed $(u_1,v_1,w_1)\in\mathcal{U} $ with $\Tr_1^k(w_1)=1$, we can consider  
		\[
		\bigl\{\, (u_1,v_1,w_1),\; (u_1+u,\;v_1+v,\;w_1+w) \, : \, (u,v,w)\in \mathcal{U}, \, \Tr_1^k(w)=1 \,\bigr\},
		\]
		as both sets span the same subspace.
		Since \(\Tr_1^k(w_1)=\Tr_1^k(w)=1\), we have \(\Tr_1^k(w_1+w)=0\).  
		In particular,  
		\((u_1+u,v_1+v,w_1+w)\in \mathcal{U}_V\), 
		where \(\mathcal{U}_V\) is defined in Equation \eqref{eq:UV}.  
		Hence, \(\mathcal{U}_V\) and \((u_1,v_1,w_1)\) generate the whole space \(\mathcal{U}\), which implies that 
		\(\dim (\mathcal{U})\leq 2k+2\). Then our assumption $n > 2k + 4$ implies that $\dim(\mathcal U) < \tfrac{n}{2} + k$,
		which yields the desired conclusion.
	\end{proof}
	
	In Proposition~\ref{thm:main}, in order to construct functions in the \(GMM\) class that do not belong to the completed Maiorana--McFarland class, we start from Maiorana--McFarland functions.
	Specifically, for
	\(\V_n^{(2)} \cong \mathbb{F}_{2^m} \times \mathbb{F}_{2^m}\) with \(n = 2m\),
	we consider the function
	\(h : \mathbb{F}_{2^m} \times \mathbb{F}_{2^m} \to \mathbb{F}_2\) defined by $h(x_1,x_2) = \Tr_{1}^{m}\!\bigl(x_1\,\pi(x_2)\bigr)$,
	where \(\pi\) is a permutation of \(\mathbb{F}_{2^m}\).
	For $a=(a_{1},a_{2}), b=(b_{1},b_{2}) \in \mathbb{F}_{2^m}\times\mathbb{F}_{2^m}$, 
	\begin{align} \label{eq:der}
		D_{a}D_{b}h(x_1,x_2) 
		&= \Tr_{1}^{m}\!\Bigl(
		x_1\,\pi(x_2)
		+\,(x_1+a_{1})\,\pi(x_2+a_{2}) \nonumber \\[4pt]
		&\quad+\,(x_1+b_{1})\,\pi(x_2+b_{2})
		+\,(x_1+a_{1}+b_{1})\,\pi(x_2+a_{2}+b_{2})
		\Bigr) \nonumber \\[6pt]
		&= \Tr_{1}^{m}\!\Bigl(
		x_1\bigl(\pi(x_2)+\pi(x_2+a_{2})+\pi(x_2+b_{2})+\pi(x_2+a_{2}+b_{2})\bigr)
		\Bigr) \nonumber \\[-1pt]
		&\quad+\;\Tr_{1}^{m}\!\Bigl(
		a_{1}\,\pi(x_2+a_{2})
		+\;b_{1}\,\pi(x_2+b_{2})
		+\;(a_{1}+b_{1})\,\pi(x_2+a_{2}+b_{2})
		\Bigr).
	\end{align}
	Recall that $ \mathcal{W} =\mathbb{F}_{2^m} \times \{0\}$
	is an $\mathcal{M}$-subspace of $h$ of dimension $m=n/2$, called the canonical $\mathcal{M}$-subspace.  
	Our goal is to determine examples of $\pi$ ensuring that $h$ has this unique $\mathcal{M}$-subspace $\mathcal{W}$ of dimension $n/2$ such that, if $\mathcal{U}$ is a nontrivial $\mathcal{M}$-subspace of $h$, then necessarily $\mathcal{U} \subseteq \mathcal{W}$. Equivalently, for distinct nonzero $a, b \in \mathbb{F}_{2^m}\times\mathbb{F}_{2^m}$,
	\[
	D_{a}D_{b}\, h(x_1,x_2)=0 
	\quad\text{for all } (x_1,x_2)\in\mathbb{F}_{2^m}\times\mathbb{F}_{2^m}
	\quad\iff\quad 
	a,b\in \mathcal{W}.
	\]
	From Equation~\eqref{eq:der}, we see that $D_{a}D_{b}\, h(x_1,x_2)=0 $ for all $(x_1,x_2)\in\mathbb{F}_{2^m}\times\mathbb{F}_{2^m}$ if and only if, for every $x_2\in \mathbb{F}_{2^m}$,
	\begin{align*}
		\pi(x_2)+\pi(x_2+a_{2})+\pi(x_2+b_{2})+\pi(x_2+a_{2}+b_{2}) &= 0, \\[4pt] 
		\Tr_{1}^{m}\!\Bigl(
		a_{1}\,\pi(x_2+a_{2})
		+\;b_{1}\,\pi(x_2+b_{2})
		+\;(a_{1}+b_{1})\,\pi(x_2+a_{2}+b_{2})
		\Bigr) &= 0.
	\end{align*}
	Hence, we obtain the following result.
	\begin{corollary}\label{cor:char}
		Let
		\(h:\mathbb{F}_{2^m}\times \mathbb{F}_{2^m}\to \mathbb{F}_{2}\)
		be an \(MM\) function defined by $h(x_1,x_2)=\Tr_{1}^{m}\!\bigl(x_1\,\pi(x_2)\bigr)$,
		where \(\pi\) is a permutation of \(\mathbb{F}_{2^m}\).
		Then $ \mathcal{W} =\mathbb{F}_{2^m} \times \{0\}$
		is the unique \(\mathcal{M}\)-subspace of \(h\) of dimension \(m\) that contains
		all nontrivial \(\mathcal{M}\)-subspaces of \(h\) if and only if the permutation
		\(\pi\) satisfies the following property.
		
		\medskip
		\noindent\textbf{Property (P).}
		For all \(x \in \mathbb{F}_{2^m}\),
		\begin{align}
			\pi(x)+\pi(x+a_{2})+\pi(x+b_{2})+\pi(x+a_{2}+b_{2})
			&= 0, \label{eq:pi1}\\[6pt]
			\Tr_{1}^{m}\!\Bigl(
			a_{1}\,\pi(x+a_{2})
			+\;b_{1}\,\pi(x+b_{2})
			+\;(a_{1}+b_{1})\,\pi(x+a_{2}+b_{2})
			\Bigr)
			&= 0, \label{eq:pi2}
		\end{align}
		if and only if one of the following holds:
		\[
		(a_1,a_2)=(0,0), \quad
		(b_1,b_2)=(0,0), \quad
		(a_1,a_2)=(b_1,b_2), \quad
		\text{or} \quad a_2=b_2=0.
		\]
	\end{corollary}
	
	\begin{remark}
		A permutation $\pi$ of $\mathbb{F}_{2^{\,m}}$ is said to satisfy
		\emph{property~(P1)} (see \cite{CaCha,PaPoKuZ}) if
		$D_a D_b \pi \neq 0$ for all linearly independent elements
		$a,b \in \mathbb{F}_{2^{\,m}}$.
		Such permutations are used to construct $MM$ functions $f(x,y)=\Tr_{1}^{m}\!\bigl(x\,\pi(y)\bigr)+h(y)$
		on $\mathbb{F}_{2^{\,m}} \times \mathbb{F}_{2^{\,m}}$ that admit no bent
		$4$-decomposition (see Section~\ref{sec:decom} for the definition) or possess a unique
		$m$-dimensional $\mathcal{M}$-subspace, namely the canonical one
		$\mathbb{F}_{2^{\,m}} \times \{0\}$.
	\end{remark}
	
	\begin{theorem}\label{thm:ex}
		Let \(m,k\) be positive integers, and let \(\pi\) be a permutation of
		\(\mathbb{F}_{2^m}\) satisfying \textbf{Property~(P)}
		of Corollary~\ref{cor:char}.
		For each \(z \in \mathbb{F}_{2^k}\), define the Boolean function
		\(f^{(z)} : \mathbb{F}_{2^m} \times \mathbb{F}_{2^m} \to \mathbb{F}_2\) by
		\[
		f^{(z)}(x_1,x_2)=
		\begin{cases}
		\Tr_1^m\!\bigl(x_1\,\pi(x_2)\bigr), & \text{if } \Tr_1^k(z)=0, \\[6pt]
		\Tr_1^m\!\bigl(x_2\,\pi(x_1)\bigr), & \text{if } \Tr_1^k(z)=1.
		\end{cases}
		\]
		If \(m > k+2\), then the function
		\[
		f :
		\mathbb{F}_{2^m} \times \mathbb{F}_{2^m} \times
		\mathbb{F}_{2^k} \times \mathbb{F}_{2^k}
		\rightarrow \mathbb{F}_2,
		\qquad
		f(x_1,x_2,y,z)
		= f^{(z)}(x_1,x_2) + \Tr_1^k(yz),
		\]
		is a \(GMM\) function that does not belong to the \(MM^{\#}\) class.
	\end{theorem}
	
	\begin{proof}
		By Corollary~\ref{cor:char}, the function $h(x_1,x_2)=\Tr_{1}^{m}\!\bigl(x_1\,\pi(x_2)\bigr)$
		admits a unique \(\mathcal{M}\)-subspace containing all nontrivial
		\(\mathcal{M}\)-subspaces of \(h\), namely $ \mathcal{W}= \mathbb{F}_{2^m}\times \{0\}$.
		Set \(\widetilde{W}=\{0\}\times \mathbb{F}_{2^m}\), and let
		\(L:\mathbb{F}_{2^m}\times\mathbb{F}_{2^m}\to
		\mathbb{F}_{2^m}\times\mathbb{F}_{2^m}\)
		be the linear permutation defined by
		\(L(x_1,x_2)=(x_2,x_1)\). Note that $\mathbb{F}_{2^m}\times\mathbb{F}_{2^m}
		=\mathcal{W} \oplus \widetilde{\mathcal{W}}$ and $L(\widetilde{\mathcal{W}} )=\mathcal{W} $.
		Then the conclusion follows directly from Proposition~\ref{thm:main}.
	\end{proof}
	
	We begin with an auxiliary result required to exhibit a permutation
	\(\pi\) satisfying \textbf{Property~(P)} of
	Corollary~\ref{cor:char}.
	
	\begin{lemma}	\label{lem:trace}
		Let $m \geq 4$ and let $c,d$ be nonzero elements of $\mathbb{F}_{2^m}$.  
		Then $\Tr_{1}^{m}\!\left(d\left(\tfrac{1}{x}+\tfrac{1}{x+c}\right)\right) $
		does not vanish identically on
		\(\mathbb{F}_{2^m}\setminus\{0,c\}\).
	\end{lemma}
	
	\begin{proof} We show that $\Tr_{1}^{m}\!\left(d\left(\tfrac{1}{x}+\tfrac{1}{x+c}\right)\right) $ is not a constant function, i.e.,
		there exist
		\(x_1,x_2\in \mathbb{F}_{2^m}\setminus\{0,c\}\)
		such that $\Tr_{1}^{m}\!\left(d\left(\tfrac{1}{x_1}+\tfrac{1}{x_1+c}\right)\right)=0$ and $\Tr_{1}^{m}\!\left(d\left(\tfrac{1}{x_2}+\tfrac{1}{x_2+c}\right)\right)=1$.
		This will establish the claim.
		
		\medskip
		
		Observe that $\Tr_{1}^{m}\!\left(d\left(\tfrac{1}{x}+\tfrac{1}{x+c}\right)\right) $ equals $0$ (respectively $1$) for some $x \in \mathbb{F}_{2^m}\setminus \{0,c\}$ if and only if the curve $\mathcal{X}_\eta$ defined by
		\[
		\mathcal{X}_\eta:\qquad
		Z^2+Z=d\left(\tfrac{1}{X}+\tfrac{1}{X+c}\right)+\eta
		\]
		has an affine $\mathbb{F}_{2^m}$-rational point for some $\eta \in \mathbb{F}_{2^m}$ with $\Tr_{1}^{m}(\eta)=0$ (respectively $\Tr_{1}^{m}(\eta)=1$).
		
		\medskip
		Let \(F_\eta\) denote the function field of \(\mathcal{X}_\eta\), namely
		\[
		F_\eta=\mathbb{F}_{2^m}(x,z),\qquad
		z^2+z=d\left(\tfrac{1}{x}+\tfrac{1}{x+c}\right)+\eta.
		\]
		Then \(F_\eta\) is an Artin--Schreier extension of
		\(\mathbb{F}_{2^m}(x)\) of degree \(2\); see \cite[Proposition~3.7.8]{sti}.
		The only ramified places of \(\mathbb{F}_{2^m}(x)\) in \(F_\eta\)
		are \((x=0)\) and \((x=c)\) (i.e., the zero of $x$ and $x+c$), each with different exponent \(2\).
		Consequently, \(\mathbb{F}_{2^m}\) is the full constant field of \(F_\eta\).
		By the Hurwitz genus formula (\cite[Theorem 3.4.13]{sti}), we have \(g(F_\eta)=1\).
		Hence, by the Hasse--Weil bound (\cite[Theorem 5.2.3]{sti}), the number \(N(F_\eta)\) of rational places of
		\(F_\eta\) satisfies
		\begin{equation}\label{eq:place}
			N(F_\eta)\geq 2^m+1-2^{(m+2)/2}.
		\end{equation}    
		
		\medskip
		Recall that an Artin--Schreier curve of the form
		\(Z^2+Z=f(X)/g(X)\) has no affine singular points whenever
		\(\gcd(f(X),g(X))=1\).
		In particular, the curve \(\mathcal{X}_\eta\) has no affine singularities.
		Since the highest-degree term in its defining equation is \(Z^2X^2\),
		there are exactly two points at infinity, namely \((0:1:0)\) and \((1:0:0)\).
		These correspond to places lying above \((x=0)\), \((x=c)\), and \((x=\infty)\), which is the pole of $x$.
		Because the places \((x=0)\) and \((x=c)\) are ramified in \(F_\eta\),
		there is a unique rational place lying above each of them.
		Moreover, depending on whether \(\Tr_{1}^{m}(\eta)=0\) or
		\(\Tr_{1}^{m}(\eta)=1\), the place \((x=\infty)\) either splits into two
		rational places or remains a single place of degree \(2\).
		Hence, there are at most four rational places corresponding to points at infinity.
		
		\medskip
		Since \(\mathcal{X}_\eta\) has no affine singular points, each affine
		\(\mathbb{F}_{2^m}\)-rational point corresponds to a unique rational place.
		Therefore, by~\eqref{eq:place}, the number \(N\) of affine rational points of
		\(\mathcal{X}_\eta\) satisfies
		\[
		N \geq 2^m-2^{(m+2)/2}-3.
		\]
		In particular, for \(m\geq 4\), there exist affine
		\(\mathbb{F}_{2^m}\)-rational points on \(\mathcal{X}_\eta\) for both
		\(\Tr_{1}^{m}(\eta)=0\) and \(\Tr_{1}^{m}(\eta)=1\),
		which yields the desired conclusion.
	\end{proof}
	
	\begin{lemma}\label{pro:pi}
		Let \(m \ge 4\), and let \(\pi\) be the permutation of \(\mathbb{F}_{2^m}\)
		defined by \(\pi(x) = x^{2^m-2}\). Then $\pi$ satisfies \textbf{Property~(P)}
		of Corollary~\ref{cor:char}. That is, $\pi(x)$
		satisfies Equations~\eqref{eq:pi1}
		and~\eqref{eq:pi2} for all $x\in\mathbb{F}_{2^m}$ if and only if one of the
		following conditions holds:
		\[
		(a_1,a_2)=(0,0), \quad (b_1,b_2)=(0,0), \quad (a_1,a_2)=(b_1,b_2),
		\quad \text{or} \quad a_2=b_2=0.
		\]
	\end{lemma}
	
	\begin{proof}
		The sufficiency of the conditions
		$(a_1,a_2)=(0,0)$, $(b_1,b_2)=(0,0)$, $(a_1,a_2)=(b_1,b_2)$, or $a_2=b_2=0$
		is straightforward. In each case, $\pi(x)$ satisfies
		Equations~\eqref{eq:pi1} and~\eqref{eq:pi2} for all $x\in\mathbb{F}_{2^m}$.
		Therefore, it remains to establish necessity.
		
		\medskip
		For any nonzero \(x\in \mathbb{F}_{2^m}\), we can write
		\(\pi(x)=1/x\). Hence, for
		\(x\in \mathbb{F}_{2^m}\setminus\{0,a_2,b_2,a_2+b_2\}\),
		Equation~\eqref{eq:pi1} can therefore be written as
		\begin{align*}
			&\pi(x)+\pi(x+a_2)+\pi(x+b_2)+\pi(x+a_2+b_2) \\[4pt]
			&= \frac{1}{x}+\frac{1}{x+a_2}+\frac{1}{x+b_2}+\frac{1}{x+a_2+b_2} \\[6pt]
			&= \frac{a_2b_2(a_2+b_2)}
			{x(x+a_2)(x+b_2)(x+a_2+b_2)}.
		\end{align*}
		Consequently, if Equation~\eqref{eq:pi1} holds for all $x\in \mathbb{F}_{2^m}$,
		then necessarily \(a_2=0\) or \(b_2=0\), or \(a_2=b_2\).
		
		\medskip
		We next consider the case \(a_2=b_2=c\) for some nonzero
		\(c\in \mathbb{F}_{2^m}\).
		For \(x\in \mathbb{F}_{2^m}\setminus\{0,c\}\),
		Equation~\eqref{eq:pi2} becomes
		\begin{align}\label{eq:a1b1}
			&\Tr_{1}^{m}\!\Bigl(
			a_{1}\,\pi(x+a_{2})
			+ b_{1}\,\pi(x+b_{2})
			+ (a_{1}+b_{1})\,\pi(x+a_{2}+b_{2})
			\Bigr) \nonumber \\[4pt]
			&= \Tr_{1}^{m}\!\Bigl(
			(a_1+b_1)\bigl(\pi(x)+\pi(x+c)\bigr)
			\Bigr) \nonumber \\[6pt]
			&= \Tr_{1}^{m}\!\Bigl(
			(a_1+b_1)\Bigl(\tfrac{1}{x}+\tfrac{1}{x+c}\Bigr)
			\Bigr).
		\end{align}
		By Lemma~\ref{lem:trace}, the above trace vanishes on
		\(\mathbb{F}_{2^m}\setminus\{0,c\}\) only if \(a_1=b_1\), i.e., $(a_1,a_2)=(b_1,b_2)$.
		
		\medskip
		Finally, suppose that exactly one of \(a_2\) or \(b_2\) is zero.
		A computation analogous to~\eqref{eq:a1b1} shows that 
		\eqref{eq:pi2} cannot hold for all \(x\in\mathbb{F}_{2^m}\) 
		unless, respectively, \(a_1=0\) or \(b_1=0\).
		
		\medskip
		This completes the proof.   
	\end{proof}
	
	As a direct consequence of Theorem~\ref{thm:ex} and Lemma~\ref{pro:pi},
	we obtain a class of functions that belong to the \(GMM\) class but not to the \(MM^{\#}\) class.
	
	\begin{corollary}\label{cor:ex1}
		Let \(m,k\) be positive integers with \(m\geq 4\) and \(m > k+2\).
		For \(z\in \mathbb{F}_{2^k}\), define the Boolean function
		\(f^{(z)}:\mathbb{F}_{2^m}\times\mathbb{F}_{2^m}\to \mathbb{F}_{2}\) by
		\[
		f^{(z)}(x_1,x_2)=
		\begin{cases}
		\Tr_1^m\!\bigl(x_1\,x_2^{2^m-2}\bigr), & \text{if } \Tr_1^k(z)=0, \\[6pt]
		\Tr_1^m\!\bigl(x_2\,x_1^{2^m-2}\bigr), & \text{if } \Tr_1^k(z)=1.
		\end{cases}
		\]
		Then the function
		\[
		f:\mathbb{F}_{2^m}\times\mathbb{F}_{2^m}\times
		\mathbb{F}_{2^k}\times\mathbb{F}_{2^k}\to \mathbb{F}_{2},
		\qquad
		f(x_1,x_2,y,z)=f^{(z)}(x_1,x_2)+\Tr_1^k(yz),
		\]
		is a \(GMM\) function which is not in the \(MM^{\#}\) class.
	\end{corollary}
	
	Similarly, we consider the Gold functions $\pi(x)=x^{2^k+1}$ introduced in \cite{gold} over the finite field $\mathbb{F}_{2^m}$. 
	It is well known that $\pi$ is a permutation of $\mathbb{F}_{2^m}$ if and only if $\gcd(2^k+1,\,2^m-1)=1$.
	Using the identity $\gcd(2^k+1,2^k-1)=1$, we obtain
	\begin{align*}
		\gcd(2^k+1,2^m-1)
		= \frac{\gcd(2^{2k}-1,2^m-1)}{\gcd(2^k-1,2^m-1)} = \frac{2^{\gcd(2k,m)}-1}{2^{\gcd(k,m)}-1}.
	\end{align*}
	Hence, $\gcd(2^k+1,2^m-1)=1$ if and only if $\gcd(k,m)=1$ and $m$ is odd. 
	
	\begin{lemma} \label{lem:gold}
		Let $m$ and $k$ be positive integers such that $\gcd(k,m)=1$ and $m$ is odd.
		Then $\pi(x)=x^{2^k+1}$ satisfies \textbf{Property~(P)}
		of Corollary~\ref{cor:char}. That is, $\pi(x)$
		satisfies Equations~\eqref{eq:pi1}
		and~\eqref{eq:pi2} for all $x\in\mathbb{F}_{2^m}$ if and only if one of the
		following conditions holds:
		\[
		(a_1,a_2)=(0,0), \quad (b_1,b_2)=(0,0), \quad (a_1,a_2)=(b_1,b_2),
		\quad \text{or} \quad a_2=b_2=0.
		\]
	\end{lemma}
	
	\begin{proof}
		Since the sufficiency of the stated conditions is immediate, we only prove necessity.
		
		A direct computation shows that
		\[
		\pi(x)+\pi(x+a_2)+\pi(x+b_2)+\pi(x+a_2+b_2)
		= a_2 b_2 \bigl(a_2^{2^k-1}+b_2^{2^k-1}\bigr).
		\]
		Since $\gcd(2^k-1,2^m-1)=1$, Equation~\eqref{eq:pi1} holds for all
		$x\in\mathbb{F}_{2^m}$ if and only if \(a_2=0\) or \(b_2=0\), or \(a_2=b_2\).
		
		Next, a straightforward computation yields
		\begin{align}\label{tr:gold}
			&\Tr_{1}^{m}\!\Bigl(
			a_{1}\,\pi(x+a_{2})
			+ b_{1}\,\pi(x+b_{2})
			+ (a_{1}+b_{1})\,\pi(x+a_{2}+b_{2})
			\Bigr) \nonumber \\[4pt]
			&\quad = \Tr_{1}^{m}\!\Bigl(
			a_1 a_2^{2^k+1}
			+ b_1 b_2^{2^k+1}
			+ (a_1+b_1)(a_2+b_2)^{2^k+1} 
			+ (a_1 b_2 + b_1 a_2)\,x^{2^k}
			+ (a_1 b_2^{2^k} + b_1 a_2^{2^k})\,x
			\Bigr)   \nonumber  \\[4pt]
			&\quad = \Tr_{1}^{m}\!\Bigl(
			a_1 a_2^{2^k+1}
			+ b_1 b_2^{2^k+1}
			+ (a_1+b_1)(a_2+b_2)^{2^k+1} 
			+ \bigl(
			(a_1 b_2 + b_1 a_2)
			+ (a_1 b_2^{2^k} + b_1 a_2^{2^k})^{2^k}
			\bigr)x^{2^k}
			\Bigr).
		\end{align}
		Therefore, Equation~\eqref{eq:pi2} holds for all $x\in\mathbb{F}_{2^m}$ only if
		the coefficient of $x^{2^k}$ vanishes, namely,
		\begin{equation}\label{eq:zero}
			(a_1 b_2 + b_1 a_2)
			+ (a_1 b_2^{2^k} + b_1 a_2^{2^k})^{2^k} = 0.
		\end{equation}
		
		We now analyze the possible cases.
		
		\medskip
		\noindent\textbf{Case (i):} $a_2=0$.
		If $b_2=0$ or $a_1=0$, then Equation~\eqref{tr:gold} vanishes trivially.
		Assume $b_2 a_1\neq 0$. Then \eqref{eq:zero} reduces to $a_1 b_2 + a_1^{2^k} b_2^{2^{2k}}=0$,
		which is equivalent to $a_1^{2^k-1} = b_2^{-(2^{2k}-1)}$.
		Since $\gcd(2^k-1,2^m-1)=1$, this holds if and only if
		$a_1=b_2^{-(2^k+1)}$. Since $m$ is odd, substituting into \eqref{tr:gold}, we obtain
		\[
		\Tr_{1}^{m}\bigl(a_1 b_2^{2^k+1}\bigr)=\Tr_{1}^{m}(1)=1,
		\]
		a contradiction. Hence, either $a_2=b_2=0$ or $(a_1,a_2)=(0,0)$.
		
		\medskip
		\noindent\textbf{Case (ii):} $b_2=0$.
		This case is symmetric to Case~(i) and implies either $a_2=b_2=0$ or $(b_1,b_2)=(0,0)$.
		
		\medskip
		\noindent\textbf{Case (iii):} $a_2=b_2=c$ for some nonzero $c\in\mathbb{F}_{2^m}$.
		Assume that $a_1\neq b_1$. Then \eqref{eq:zero} becomes
		$c(a_1+b_1) + c^{2^{2k}}(a_1+b_1)^{2^k}=0$,
		which is equivalent to
		$\bigl(c^{2^k+1}(a_1+b_1)\bigr)^{2^k-1}=1$.
		Since $\gcd(2^k-1,2^m-1)=1$, this holds if and only if
		$c^{2^k+1}(a_1+b_1)=1$. Substituting into \eqref{tr:gold}, we obtain
		\[
		\Tr_{1}^{m}\bigl(c^{2^k+1}(a_1+b_1)\bigr)
		= \Tr_{1}^{m}(1)=1,
		\]
		a contradiction. Therefore, $a_1=b_1$, and hence $(a_1,a_2)=(b_1,b_2)$.
		
		\medskip
		Combining all cases completes the proof.
	\end{proof}
	
	Similarly, Theorem~\ref{thm:ex} together with Lemma~\ref{lem:gold} yields another
	class of functions that belong to the \(GMM\) class but not to the completed
	\(MM\) class.
	
	
	\begin{corollary}\label{cor:ex2}
		Let $m$ and $k$ be positive integers such that $\gcd(k,m)=1$ and $m$ is odd with \(m > k+2\).
		For $z\in\mathbb{F}_{2^k}$, define the Boolean function
		$f^{(z)}:\mathbb{F}_{2^m}\times\mathbb{F}_{2^m}\to\mathbb{F}_{2}$ by
		\[
		f^{(z)}(x_1,x_2)=
		\begin{cases}
		\Tr_1^m\!\bigl(x_1\,x_2^{2^k+1}\bigr), & \text{if } \Tr_1^k(z)=0, \\[6pt]
		\Tr_1^m\!\bigl(x_2\,x_1^{2^k+1}\bigr), & \text{if } \Tr_1^k(z)=1.
		\end{cases}
		\]
		Then the function
		\[
		f:\mathbb{F}_{2^m}\times\mathbb{F}_{2^m}\times
		\mathbb{F}_{2^k}\times\mathbb{F}_{2^k}\to\mathbb{F}_{2},
		\qquad
		f(x_1,x_2,y,z)=f^{(z)}(x_1,x_2)+\Tr_1^k(yz),
		\]
		is a \(GMM\) function that is not in the \(MM^{\#}\) class.
	\end{corollary}
	
	\begin{remark}
		We recall that any Boolean bent function on $\V_n^{(2)}$ belonging to the completed $\mathcal{PS}_{ap}$ class has algebraic degree $n/2$. Consequently, the bent functions in Corollaries~\ref{cor:ex1} and~\ref{cor:ex2} lie neither in the \(MM^{\#}\) class nor in the $\mathcal{PS}_{ap}^{\#}$ class.
	\end{remark}
	
	
	\section{Decomposition of the generalized $\mathcal{PS}_{ap}$ functions } \label{sec:decom}
	
	Let $\mathcal{S} \subseteq \V_n^{(2)}$ be a subspace of dimension $n-2$, and
	let $\mathcal{W} \subseteq \V_n^{(2)}$ be a complementary subspace such that $\V_n^{(2)} = \mathcal{S} \oplus \mathcal{W}$.
	Let $f \colon \V_n^{(2)} \to \mathbb{F}_2$ be a Boolean function.
	For each $w_i \in \mathcal{W}$, define $f_i(x) = f(x + w_i)$ for $x \in \mathcal{S}$.
	Equivalently, $f_i$ is the restriction of $f$ to the coset
	$w_i + \mathcal{S}$, viewed as a Boolean function on $\mathcal{S}$.
	
	The $4$-decomposition of $f$ with respect to $\mathcal{S}$ is then defined as the
	sequence $(f_1,f_2,f_3,f_4)$,
	where $\{w_1,w_2,w_3,w_4\} = \mathcal{W}$. Throughout the paper, we simply refer to this as a decomposition (with
	respect to $\mathcal{S}$).
	If all the functions $f_i$ are bent (respectively, semibent), then $f$
	is said to admit a bent (respectively, semibent) decomposition with
	respect to $\mathcal{S}$. 
	
	We remark that, in the decomposition of a function $f$, the
	functions $f_i$ are pairwise EA-equivalent and hence have
	the same extended Walsh spectrum. In particular, they are bent or semibent
	simultaneously. Therefore, it suffices to consider the restriction of
	$f$ to the subspace $\mathcal{S}$.
	Moreover, since $\mathcal S$ has codimension $2$, there exist linearly independent vectors 
	$u,v \in \V_n^{(2)}$ such that $\mathcal S = \langle u,v\rangle^\perp$.
	In other words,
	\[
	\mathcal S
	=
	\{\, x \in \V_n^{(2)} : \langle u,x\rangle_n = \langle v,x\rangle_n = 0 \,\}.
	\]
	An equivalent criterion for a bent function $f$ to admit a bent
	(respectively, semibent) decomposition, expressed in terms of the
	second-order derivative of its dual function, is given in
	\cite[Theorem~7]{CaCha} as follows.
	
	\begin{lemma}
		\label{lem:decom}
		Let $f \colon \V_n^{(2)} \to \F_2$ be a bent function with dual $f^*$.
		Let $u,v \in \V_n^{(2)}$ be linearly independent, and let $\mathcal{S} = \langle u,v \rangle^\perp$
		be the orthogonal complement of $\langle u,v \rangle$.
		Then the following hold:
		\begin{itemize}
			\item[(i)] $f$ admits a bent decomposition with respect to $\mathcal{S}$
			if and only if $D_u D_v f^* = 1$.
			\item[(ii)] $f$ admits a semibent decomposition with respect to $\mathcal{S}$
			if and only if $D_u D_v f^* = 0$.
		\end{itemize}
	\end{lemma}
	
	From Lemma~\ref{lem:decom}, it follows that $f$ admits a semibent
	decomposition if and only if its dual $f^*$ possesses a nontrivial
	$\mathcal{M}$-subspace; equivalently, its linearity index satisfies
	$\mathrm{ind}(f^*) > 1$. In particular, any $MM$ or $GMM$ bent function
	admits a semibent decomposition. In contrast to $MM$ and $GMM$ bent
	functions, it appears that the majority of the $\mathcal{PS}_{ap}$ bent
	functions admit neither a bent nor a semibent decomposition; for
	details, we refer to \cite{akkmpp26}.
	
	In this section, we investigate the decomposition of bent functions in the
	generalized $\mathcal{PS}_{ap}$ class. 
	As in~\cite{akkmpp26}, our approach relies on curves defined over finite fields 
	and on estimates for their number of affine rational points. 
	However, due to the definition of generalized $\mathcal{PS}_{ap}$ functions, 
	the analysis requires more involved technical computations.
	
	Let \(m\), \(k\), and \(e\) be integers such that \(k \mid m\),
	\(e \equiv 2^{\ell} \mod{(2^{k}-1)}\), and
	\(\gcd(2^{m}-1, e) = 1\).
	Let \(\eta\) denote the multiplicative inverse of \(e\) modulo \((2^{m}-1)\), that
	is, $\eta e \equiv 1 \mod{(2^{m}-1)}$.
	In particular, we consider the generalized $\mathcal{PS}_{ap}$ bent functions defined by
	\begin{equation}\label{eq:gPSap}
		g(x,y)=
		P\!\left(\Tr_{k}^{m}\!\left(x\, y^{-\eta}\right)\right),
	\end{equation}
	where $P \colon \mathbb{F}_{2^{k}} \to \mathbb{F}_2$ is a balanced Boolean
	function.
	
	To apply Lemma~\ref{lem:decom}, we first determine the dual of the function \(g\)
	defined in \eqref{eq:gPSap}.
	
	\begin{lemma}
		Let \(g \colon \mathbb{F}_{2^{m}} \times \mathbb{F}_{2^{m}} \to \mathbb{F}_2\)
		be the generalized \(\mathcal{PS}_{ap}\) bent function defined in
		\eqref{eq:gPSap}. Then the dual bent function \(g^{\ast}\) is given by
		\[
		g^{\ast}(x,y)
		=
		P\!\left(\Tr^{m}_{k}\!\left(y\, x^{-e}\right)^{2^{m-\ell}}\right)
		=
		P\!\left(\Tr^{m}_{k}\!\left(\tilde y\, \tilde x^{-e}\right)\right),
		\]
		where $\tilde y=y^{2^{m-\ell}}$ and $\tilde x=x^{2^{m-\ell}}$.
	\end{lemma}
	\begin{proof}
		For \((u,v) \in \bigl(\mathbb{F}_{2^{m}} \times \mathbb{F}_{2^{m}}\bigr)
		\setminus \{(0,0)\}\), let \(\chi_{u,v}\) denote the character of
		\(\mathbb{F}_{2^{m}} \times \mathbb{F}_{2^{m}}\) defined by
		$\chi_{u,v}(x,y)
		=
		(-1)^{\Tr_1^{m}(ux + vy)}$.
		Let $\Gamma_2 = \{V, \mathcal{B}(\gamma)\; : \; \gamma\in\F_{2^k}\}$ 
		be the generalized Desarguesian spread defined in \eqref{eq:Gamma}, where \(V\) and
		\(\mathcal{B}(\gamma)\)  are as in \eqref{eq:Bgamma} and \eqref{eq:Agamma}, respectively.
		We note that \(g(x,y)=P(\gamma)\) for all \((x,y)\in \mathcal{B}(\gamma)\) and
		\(g(x,y)=P(0)\) for all \((x,y)\in V\), i.e., \(g\) is a bent function arising from the bent partition
		\(\Gamma_2\).
		By~\cite[Proposition~10]{akm24}, for \(\gamma \in \mathbb{F}_{2^{k}}\) we have
		\[
		\chi_{u,v}\bigl(B(\gamma)\bigr)
		=\sum_{(x,y)\in B(\gamma)}(-1)^{\Tr_1^{m}(ux + vy)}=
		\begin{cases}
		2^{m} - 2^{m-k}, &
		\text{if } u \neq 0 \text{ and }
		\gamma^{2^\ell} = \Tr_k^{m}\!\left(vu^{-e}\right), \\[1mm]
		-\,2^{m-k}, & \text{otherwise},
		\end{cases}
		\] 
		and
		\[
		\chi_{u,v}(V)
		=\sum_{(x,y)\in V}(-1)^{\Tr_1^{m}(ux + vy)}
		=
		\begin{cases}
		0, & \text{if } u \neq 0, \\[1mm]
		2^{m}, & \text{otherwise}.
		\end{cases}
		\]
		We now compute the Walsh transform of \(g\):
		\begin{align*}
			W_g(u,v)
			&=
			\sum_{(x,y) \in \mathbb{F}_{2^{m}} \times \mathbb{F}_{2^{m}}}
			(-1)^{\,g(x,y) + \Tr_1^{m}(ux+vy)} \\
			&=
			\sum_{(x,y) \in \mathbb{F}_{2^{m}} \times \mathbb{F}_{2^{m}}}
			(-1)^{\,P\!\left(\Tr_{k}^{m}\!\left(x\,y^{-\eta }\right)\right)
				+ \Tr_1^{m}(ux+vy)} \\
			&=
			\sum_{\gamma \in \mathbb{F}_{2^{k}}}
			\sum_{(x,y) \in B(\gamma)}
			(-1)^{\,P(\gamma) + \Tr_1^{m}(ux+vy)}
			+
			\sum_{(x,y) \in V}
			(-1)^{\,P(0) + \Tr_1^{m}(ux+vy)} \\
			&=
			\sum_{\gamma \in \mathbb{F}_{2^{k}}}
			(-1)^{\,P(\gamma)} \chi_{u,v}\bigl(B(\gamma)\bigr)
			+
			(-1)^{\,P(0)} \chi_{u,v}(V).
		\end{align*}
		Using the character values above together with the balancedness of
		\(P \colon \mathbb{F}_{2^{k}} \to \mathbb{F}_2\), we obtain
		\[
		W_g(u,v)
		=
		\begin{cases}
		(-1)^{P(0)}\,2^{m}, & \text{if } u = 0, \\[1mm]
		(-1)^{P(\gamma)}\,2^{m},
		& \text{if } u \neq 0 \text{, where }
		\gamma^{2^\ell} = \Tr_k^{m}\!\left(v\,u^{-e}\right).
		\end{cases}
		\]
		This yields the claimed expression for the dual function, together with the identity $\gamma = \gamma^{2^m}
		= \Tr_k^{m}\!\left(v\,u^{-e}\right)^{\,2^{m-\ell}}
		=\Tr_k^{m}\!\left(v^{\,2^{m-\ell}}\,u^{-e{\,2^{m-\ell}}}\right)$.
	\end{proof}
	
	Note that the trace map
	\(\Tr_k^{m} \colon \mathbb{F}_{2^{m}} \to \mathbb{F}_{2^{k}}\)
	and the function
	\(P \colon \mathbb{F}_{2^{k}} \to \mathbb{F}_2\)
	are both balanced. Consequently, the composition
	\(P\!\left(\Tr_k^{m}(z)\right)\)
	defines a balanced Boolean function from \(\mathbb{F}_{2^{m}}\) to \(\mathbb{F}_2\).
	Hence, there exists a permutation \(Q\) of \(\mathbb{F}_{2^{m}}\) such that $P\!\left(\Tr_k^{m}(z)\right)
	=
	\Tr_1^{m}\!\left(Q(z)\right)$.
	Therefore, in the remainder of this section, we consider the generalized
	\(\mathcal{PS}_{ap}\) bent functions from
	\(\mathbb{F}_{2^{m}} \times \mathbb{F}_{2^{m}}\) to \(\mathbb{F}_2\)
	of the form 
	\begin{align}\label{gps:new}
		f(x,y)
		=
		\Tr_1^{m}\!\left(Q\!\left(x\,y^{-\eta }\right)\right).
	\end{align}
	We remark that, in this case, the dual function of \(f\) is given by
	\begin{align}\label{gps:newdual}
		f^{\ast}(x,y)
		=
		\Tr_1^{m}\!\left(Q\!\left(\tilde y\, \tilde x^{-e}\right)\right),
	\end{align}
	where $\tilde y=y^{2^{m-\ell}}$ and $\tilde x=x^{2^{m-\ell}}$.
	Moreover, by Lemma~\ref{lem:second}, we can without loss of generality suppose that $Q(0)=0$.
	
	
	\begin{proposition}\label{mainprop}
		Let $m > 4$, and let $Q$ be a permutation of $\mathbb{F}_{2^{m}}$ of odd polynomial degree less than or equal to $\frac{2^{m/4}-1}{3e}$. 
		Let $u=(a,b)$ and $v=(c,d)$ be two linearly independent vectors in $\mathbb{F}_{2^{m}} \times \mathbb{F}_{2^{m}}$, and set $\mathcal{S} = \langle u, v \rangle^{\perp}$. 
		Let \(f \colon \mathbb{F}_{2^{m}} \times \mathbb{F}_{2^{m}} \to \mathbb{F}_2\) be the generalized
		\(\mathcal{PS}_{ap}\) bent function defined in \eqref{gps:new}.
		If $ad + bc \neq 0$, then $f$ admits neither a bent nor a semibent decomposition on $\mathcal{S}$.
	\end{proposition}
	
	
	\begin{remark} \label{rem:mainpro}
		Let
		$f\colon \mathbb{F}_{2^{m}} \times \mathbb{F}_{2^{m}} \to \mathbb{F}_2$
		be the function defined in \eqref{gps:new}, and suppose that it satisfies the
		assumptions of Proposition~\ref{mainprop}. Then, equivalently, we have:
		\begin{itemize}
			\item[(i)] The dual function $f^{\ast}$ of $f$ does not admit any
			two-dimensional $\mathcal{M}$-subspace
			$\langle (a,b),(c,d) \rangle$ with $ad+bc \neq 0$.
			
			\item[(ii)] If $f$ admits a bent or semibent decomposition on
			$\langle (a,b),(c,d) \rangle^{\perp}$, then necessarily $ad+bc=0$.
			Equivalently, $(c,d)=\lambda (a,b)$ for some
			$\lambda \in \mathbb{F}_{2^{m}} \setminus \mathbb{F}_{2}$.
		\end{itemize}
	\end{remark}
	
	Before proving Proposition~\ref{mainprop}, we require some preliminary results.
	The first lemma is analogous to the one appearing in~\cite{akkmpp26}; however, 
	we include its proof for completeness.
	
	\begin{lemma}\label{ftof}
		Let $\mathcal{S}$ be the $(n-2)$-dimensional subspace of $\F_{2^m}\times \F_{2^m}$ defined by
		\[
		\Tr_1^m(ax+by)=0
		\quad\text{and}\quad
		\Tr_1^m(cx+dy)=0.
		\]
		Assume that $ad+bc\neq 0$. Then the function $f(x,y)=\Tr_1^m\!\bigl(Q(xy^{-\eta})\bigr)$
		admits a bent decomposition (respectively, a semibent decomposition) on $\mathcal{S}$ if and only if $D_{(1,0)}D_{(0,1)}\widehat{f}(x,y)$
		is identically equal to $1$ (respectively, identically equal to $0$) for all \((x,y) \in \F_{2^m}\times \F_{2^m}\), where
		\[
		\widehat{f}(x,y)
		=
		\Tr_1^m\!\left(
		Q\!\left(
		\bigl(b^{2^{m-\ell}} x + d^{2^{m-\ell}} y\bigr)
		\bigl(a^{2^{m-\ell}} x + c^{2^{m-\ell}} y\bigr)^{-e}
		\right)
		\right).
		\]
	\end{lemma}
	\begin{proof}
		By Lemma~\ref{lem:decom}, the function $f$ admits a bent decomposition (respectively, a semibent decomposition) on $\mathcal{S}$ if and only if the second-order derivative $D_u D_v f^*(x)$
		of the dual function $f^*(x,y)$ given in \eqref{gps:newdual}
		is constantly equal to $1$ (respectively, constantly equal to $0$), where $u=(a,b)$ and $v=(c,d)$.
		
		By Lemma~\ref{lem:second}, this condition is equivalent to requiring that $D_{(1,0)} D_{(0,1)} f^*\bigl(L(x,y)\bigr)$
		is constantly equal to $1$ (respectively, $0$), where $L$ is a linear permutation of
		$\F_{2^m}\times\F_{2^m}$ satisfying $L(1,0)=(a,b)$ and $L(0,1)=(c,d)$.
		
		Define the linear map $L$ by $L(x,y)=(ax^{2^\ell}+cy^{2^\ell},\; bx^{2^\ell}+dy^{2^\ell})$.
		The map $L$ permutes $\F_{2^m}\times\F_{2^m}$ if and only if $ad+bc\neq 0$. Under this assumption, we indeed have
		$L(1,0)=(a,b)$ and $L(0,1)=(c,d)$. Moreover,
		\[
		\begin{aligned}
		f^*\bigl(L(x,y)\bigr)
		&=
		\Tr_1^m\!\left(
		Q\!\left(
		(bx^{2^\ell}+dy^{2^\ell})^{2^{m-\ell}}
		(ax^{2^\ell}+cy^{2^\ell})^{-2^{m-\ell}e}
		\right)
		\right) \\[4pt]
		&=
		\Tr_1^m\!\left(
		Q\!\left(
		\bigl(b^{2^{m-\ell}}x+d^{2^{m-\ell}}y\bigr)
		\bigl(a^{2^{m-\ell}}x+c^{2^{m-\ell}}y\bigr)^{-e}
		\right)
		\right) \\
		&= \widehat{f}(x,y).
		\end{aligned}
		\]
		This yields the desired conclusion.
	\end{proof}

	By Lemma~\ref{ftof}, in order to prove Proposition~\ref{mainprop}, it suffices to show that,
	for $u=(1,0)$ and $v=(0,1)$, the second-order derivative $D_uD_v\widehat f(x,y)$
	is nonconstant.
	Since $a^{2^{m-\ell}} d^{2^{m-\ell}}
	+
	b^{2^{m-\ell}} c^{2^{m-\ell}}
	=
	(ad+bc)^{2^{m-\ell}}$,
	the condition $ad+bc \neq 0$ holds if and only if $a^{2^{m-\ell}} d^{2^{m-\ell}}
	+
	b^{2^{m-\ell}} c^{2^{m-\ell}}
	\neq 0$.
	Hence, without loss of generality, we may assume $ad+bc \neq 0$ and consider
	\[
	\widetilde{f}(x,y)
	=
	\Tr_1^m\!\left(
	Q\!\left(
	(bx+dy)(ax+cy)^{-e}
	\right)
	\right).
	\]
	More precisely, we compute
	\begin{align}
		D_uD_v\widetilde f(x,y)
		&=
		\widetilde f(x+1,y+1)
		+\widetilde f(x+1,y)
		+\widetilde f(x,y+1)
		+\widetilde f(x,y) \label{eq:Duv} \\
		&=
		\Tr_1^m\Big(
		Q\big((b(x+1)+d(y+1))(a(x+1)+c(y+1))^{-e}\big) \nonumber \\
		&\qquad
		+ Q\big((b(x+1)+dy)(a(x+1)+cy)^{-e}\big) \nonumber \\
		&\qquad
		+ Q\big((bx+d(y+1))(ax+c(y+1))^{-e}\big) \nonumber \\
		&\qquad
		+ Q\big((bx+dy)(ax+cy)^{-e}\big)
		\Big). \nonumber
	\end{align}
	
	The following lemma is frequently used to characterize the circumstances under which the second-order derivative of a function cannot be constant.
	
	\begin{lemma}\label{lem:curveDuv}
		For a positive integer $\ell \geq 1$, let $\mathcal{X}_\vartheta$ be the curve over $\F_{2^m}$ defined by
		\begin{equation} \label{eq:X_theta}
			Z^2+Z
			=
			P\!\left( \frac{\alpha_1 X+\beta_1}{\kappa_1 X+\upsilon_1} \right)
			+
			\cdots
			+
			P\!\left( \frac{\alpha_\ell X+\beta_\ell}{\kappa_\ell X+\upsilon_\ell} \right)
			+\vartheta ,
		\end{equation}
		where $P$ is a polynomial of odd degree $t$.
		Assume that
		\[
		\alpha_i\upsilon_i+\beta_i\kappa_i \neq 0
		\quad \text{for } i=1,\ldots, \ell,
		\qquad\text{and}\qquad
		\exists j \text{ s.t. }\frac{\upsilon_j}{\kappa_j} \neq \frac{\upsilon_i}{\kappa_i}
		\ \text{for all } i\neq j.
		\]
		Then the number $N_\vartheta$ of affine $\F_{2^m}$-rational points
		$(x,z)\in \mathcal{X}_\vartheta$ with  $x\neq \upsilon_i/\kappa_i$, for $i=1, \ldots, \ell$, of $\mathcal{X}_\vartheta$
		satisfies
		\[
		N_\vartheta
		\;\ge\;
		2^m - 1 - (\ell t+1)\ell t\,2^{\frac{m}{2}} - \ell (\ell t+2).
		\]
		In particular, if $m \ge 4$ and $\ell t \le 2^{m/4}-1$, then $N_\vartheta > 0$.
	\end{lemma}
	
	\begin{proof}
		Let $F_\vartheta=\F_{2^m}(x,z)$ be the function field of $\mathcal{X}_\vartheta $.
		We may regard $F_\vartheta$ as an extension of the rational function field
		$\F_{2^m}(x)$ defined by the equation $z^2+z=g_\vartheta(x)$,
		where
		\[
		g_\vartheta(x)
		=
		P\!\left( \frac{\alpha_1 x+\beta_1}{\kappa_1 x+\upsilon_1} \right)
		+
		\cdots
		+
		P\!\left( \frac{\alpha_\ell x+\beta_\ell}{\kappa_\ell x+\upsilon_\ell} \right)
		+\vartheta.
		\]
		Note that $F_\vartheta/\F_{2^m}(x)$ is an Artin--Schreier extension of degree~$2$;
		see \cite[Proposition~3.7.8]{sti}.
		A place $R$ of $\F_{2^m}(x)$ ramifies in $F_\vartheta$ only if the valuation
		$v_R(g_\vartheta(x))$ is negative.
		This can occur only when $R=(x=\upsilon_i/\kappa_i)$ for $i=1, \ldots, \ell$.
		Moreover, by the strict triangle inequality (see \cite[Lemma~1.1.11]{sti}), 
		for the place $R_j$ corresponding to 
		$x = \upsilon_j/\kappa_j$, where $\tfrac{\upsilon_j}{\kappa_j} \neq \tfrac{\upsilon_i}{\kappa_i}$ for all $i \neq j$,
		we obtain
		\[
		v_{R_j}\bigl(g_\vartheta(x)\bigr) = -\deg(P) = -t.
		\]
		Since $t$ is odd, the place $R_j$ is totally ramified in the extension
		$F_\vartheta/\F_{2^m}(x)$.
		Consequently, $F_\vartheta$ is a function field with full constant field
		$\F_{2^m}$.
		It follows that $\mathcal{X}_\vartheta$ is an absolutely irreducible curve defined
		over $\F_{2^m}$; see \cite[Corollary~3.6.8]{sti}.
		Furthermore, $\mathcal{X}_\vartheta $ has degree
		$\deg(\mathcal{X}_\vartheta)\le \ell t+2$.
		Applying the Hasse-Weil bound (see \cite[Theorem~9.57]{hkt}), we obtain the following
		estimate for the number $N(\mathcal{X}_\vartheta)$ of $\F_{2^m}$-rational points of
		$\mathcal{X}_\vartheta$ in the projective plane:
		\begin{align}\label{eq:N}
			N(\mathcal{X}_\vartheta)
			&\ge
			2^m+1-(\deg(\mathcal{X}_\vartheta)-1)(\deg(\mathcal{X}_\vartheta)-2)\,2^{\frac{m}{2}} \nonumber\\
			&\ge
			2^m+1-(\ell t+1)\ell t \,2^{\frac{m}{2}}.
		\end{align}
		As the highest-degree term in the defining equation of
		$\mathcal{X}_\vartheta$ is $X^{s}Z^2$ with $s\le \ell t$, the curve
		$\mathcal{X}_\vartheta$ has at most two $\F_{2^m}$-rational points at infinity, namely $(0:1:0)$ and $(1:0:0)$.
		In order to ensure that the denominators in \eqref{eq:X_theta} do not vanish, 
		we must also exclude all points lying on the $\ell$ lines defined by $\kappa_i X + \upsilon_i = 0$.
		By Bézout’s theorem, a line intersects $\mathcal{X}_\vartheta$ in at most $\deg(\mathcal{X}_\vartheta)$ points.
		Hence, by subtracting $2+\ell\deg(\mathcal{X}_\vartheta)$ from \eqref{eq:N}, we obtain the desired bound on the number $N_\vartheta$ of affine $\F_{2^m}$-rational points. 
		Moreover, whenever $m \ge 4$
		and $\ell t \le 2^{m/4}-1$, we obtain
		\[
		(\ell t+1)\ell t\,2^{\frac{m}{2}} + \ell(\ell t+2) + 1 < 2^{m},
		\]
		and hence $N_\vartheta > 0$.
	\end{proof}
	
	\medskip
	
	\noindent {\it \textbf{Proof of Proposition \ref{mainprop}.}}
	The proof is carried out by a case-by-case analysis.
	
	\medskip
	
	\noindent \textbf{Case (i):} $b=d$.  \\[.3em]
	Note that the condition $b=d$ implies $b\neq 0$ and $a\neq c$, since
	$ad+bc\neq 0$.
	In this case, we set $y=x$.
	Then, by \eqref{eq:Duv}, for $x\neq \tfrac{a}{a+c},\ \tfrac{c}{a+c}$,
	we obtain
	\begin{align*}
		D_uD_v\widetilde{f}(x,x)
		&=
		\Tr^m_1\!\left(
		Q\bigl(b((a+c)x+a)^{-e}\bigr)
		+
		Q\bigl(d((a+c)x+c)^{-e}\bigr)
		\right) \\
		&=
		\Tr^m_1\!\left(
		Q\!\left( \frac{b}{((a+c)x+a)^{e}} \right)
		+
		Q\!\left( \frac{b}{((a+c)x+c)^{e}} \right)
		\right).
	\end{align*}
	Since $\gcd(e,2^m-1)=1$, the map $x\mapsto x^e$ permutes $\F_{2^m}$.
	Thus, there exists $\tilde b\in\F_{2^m}$ such that $b=\tilde b^{\,e}$.
	Setting $P(X)=Q(X^e)$, we see that there exists
	$x\in \F_{2^m}\setminus\left\{\frac{a}{a+c},\frac{c}{a+c}\right\}$
	such that
	$D_uD_v\tilde f(x,x)=0$, respectively $D_uD_v\tilde f(x,x)=1$,
	if and only if the curve
	\[
	Z^2+Z
	=
	P\!\left( \frac{\tilde b}{(a+c)X+a} \right)
	+
	P\!\left( \frac{\tilde b}{(a+c)X+c} \right)
	+\vartheta
	\]
	has an affine rational point $(x,z)$ with
	$x\notin \left\{\frac{a}{a+c},\frac{c}{a+c}\right\}$,
	where $\vartheta\in\F_{2^m}$ satisfies
	$\Tr^m_1(\vartheta)=0$, respectively $\Tr^m_1(\vartheta)=1$.
	
	By setting in \eqref{eq:X_theta} $\ell=2$,
	\[
	(\alpha_1,\beta_1,\kappa_1,\upsilon_1)=(0,\tilde b,a+c,a),
	\quad \text{and} \quad
	(\alpha_2,\beta_2,\kappa_2,\upsilon_2)=(0,\tilde b,a+c,c),
	\]
	this condition is equivalent to the curve $\mathcal{X}_\vartheta$
	having an affine $\F_{2^m}$-rational point $(x,z)$ with
	$x\notin \left\{\frac{a}{a+c},\frac{c}{a+c}\right\}$.
	Note that $\alpha_1\upsilon_1+\beta_1\kappa_1\neq 0$ and $\alpha_2\upsilon_2+\beta_2\kappa_2\neq 0$
	since $b\neq 0$ and $a\neq c$.
	Moreover, $\tfrac{\upsilon_1}{\kappa_1}\neq \tfrac{\upsilon_2}{\kappa_2}$,
	as $a\neq c$.
	Therefore, the existence of such a point follows from
	Lemma~\ref{lem:curveDuv}.
	
	\medskip
	
	\noindent \textbf{Case (ii):} $bd\neq 0$ and $b\ne d$.  \\[.3em]
	Set $y=\tfrac{b}{d}x$ and $\delta=a+\tfrac{cb}{d}$.
	By the assumption $ad+bc\neq 0$, we have $\delta\neq 0$.
	Then, for $x\neq \tfrac{a+c}{\delta},\ \tfrac{a}{\delta},\ \tfrac{c}{\delta}$,
	we obtain
	\begin{align}\label{eq:case2}
		D_uD_v\widetilde{f}\!\left(x,\frac{b}{d}x\right)
		&=
		\Tr^m_1\!\left(
		Q\!\left( \frac{b+d}{(\delta x+(a+c))^{e}} \right)
		+
		Q\!\left( \frac{b}{(\delta x+a)^{e}} \right)
		+
		Q\!\left( \frac{d}{(\delta x+c)^{e}} \right)
		\right).
	\end{align}
	As the map $x\mapsto x^e$ permutes $\F_{2^m}$,
	there exist elements
	$\tilde b,\tilde d,\widetilde{b+d}\in\F_{2^m}$ such that
	\[
	b=\tilde b^{\,e},\qquad
	d=\tilde d^{\,e},\qquad
	b+d=\left( \widetilde{b+d}\right)^{\,e}.
	\]
	Defining $P(X)=Q(X^e)$, Equation \eqref{eq:case2} can be rewritten as
	\[
	D_uD_v\widetilde{f}\!\left(x,\frac{b}{d}x\right)
	=
	\Tr^m_1\!\left(
	P\!\left( \frac{\widetilde{b+d}}{\delta x+(a+c)} \right)
	+
	P\!\left( \frac{\tilde b}{\delta x+a} \right)
	+
	P\!\left( \frac{\tilde d}{\delta x+c} \right)
	\right).
	\]
	Consequently, there exists $x\in \F_{2^m}\setminus
	\left\{\tfrac{a+c}{\delta},\tfrac{a}{\delta},\tfrac{c}{\delta}\right\}$
	such that
	$D_uD_v\tilde{f}(x,\frac{b}{d}x)=0$,
	respectively $D_uD_v\tilde{f}(x,\frac{b}{d}x)=1$,
	if and only if the curve defined by
	\[
	Z^2+Z
	=
	P\!\left( \frac{\widetilde{b+d}}{\delta X+(a+c)} \right)
	+
	P\!\left( \frac{\tilde b}{\delta X+a} \right)
	+
	P\!\left( \frac{\tilde d}{\delta X+c} \right)
	+\vartheta
	\]
	has an affine $\F_{2^m}$-rational point $(x,z)$ with $x\notin
	\left\{\tfrac{a+c}{\delta},\tfrac{a}{\delta},\tfrac{c}{\delta}\right\}$,
	where $\vartheta\in\F_{2^m}$ satisfies
	$\Tr^m_1(\vartheta)=0$, respectively $\Tr^m_1(\vartheta)=1$.
	
	
	We now verify that the assumptions of \eqref{lem:curveDuv} are satisfied. 
	To this end, in \eqref{eq:X_theta} we set $\ell=3$ and 
	\[
	(\alpha_1,\beta_1,\kappa_1,\upsilon_1)
	=
	(0,\widetilde{b+d},\delta,a+c), \,
	(\alpha_2,\beta_2,\kappa_2,\upsilon_2)
	=
	(0,\tilde b,\delta,a),
	\,
	(\alpha_3,\beta_3,\kappa_3,\upsilon_3)
	=
	(0,\tilde d,\delta,c).
	\]
	Note that for every $i$ we have $\alpha_i \upsilon_i + \beta_i \kappa_i
	=
	\beta_i \delta$,
	which is nonzero since $\delta \neq 0$, $bd \neq 0$, and $b \neq d$.
	Moreover, as the $\kappa_i$ are identical, the second assumption reduces to showing that at least one of the $\upsilon_i$ differs from the other two. 
	The equalities $a+c = a = c$ would hold only if $a=0$ and $c=0$, which is impossible because $ad + bc \neq 0$.
	Therefore, by \eqref{lem:curveDuv}, such a rational point always exists.
	
	\medskip
	
	\noindent \textbf{Case (iii):} $b=0$ or $d=0 $. \\[.3em]
	Note that $b$ and $d$ cannot vanish simultaneously since $ad+bc \neq 0$.
	As the case $d=0$ is analogous to the case $b=0$, we may assume without loss of generality that $b=0$.
	Then $ad \neq 0$, and hence
	\begin{align*}
		D_uD_v\widetilde{f}(x,y)
		&= \Tr^m_1\Big(
		Q\big(d(y+1)(a(x+1)+c(y+1))^{-e}\big)
		+Q\big(dy(a(x+1)+cy)^{-e}\big)  \\
		&\qquad\quad
		+Q\big(d(y+1)(ax+c(y+1))^{-e}\big)
		+Q\big(dy(ax+cy)^{-e}\big)
		\Big).
	\end{align*}
	We now set $y=0$. Then, for $x \neq  \tfrac{c}{a},\, 1+\tfrac{c}{a}$, we obtain
	\begin{align*}
		D_uD_v\widetilde{f}(x,0)
		= \Tr^m_1\left(
		Q\left( \frac{d}{(ax+(a+c))^{e}}\right)
		+Q\left( \frac{d}{(ax+c)^{e}}\right)
		\right).
	\end{align*}
	Consequently, there exists
	$x \in \mathbb{F}_{2^m} \setminus \{ \tfrac{c}{a},\, 1+\tfrac{c}{a} \}$
	such that $D_uD_v\tilde{f}(x,0)=0$, respectively $D_uD_v\tilde{f}(x,0)=1$,
	if and only if the curve defined by
	\begin{align*}
		Z^2+Z
		=
		P\left( \frac{\tilde{d}}{aX+(a+c)} \right)
		+P\left( \frac{\tilde{d}}{aX+c} \right)
		+\vartheta
	\end{align*}
	has an affine $\F_{2^m}$-rational point $(x,z)$ with 
	$x \notin \{ \tfrac{c}{a},\, 1+\tfrac{c}{a} \}$,
	where $P(X)=Q(X^e)$, $d=\tilde{d}^e$, and
	$\vartheta \in \mathbb{F}_{2^m}$ satisfies
	$\Tr^m_1(\vartheta)=0$, respectively $\Tr^m_1(\vartheta)=1$.
	By setting in \eqref{eq:X_theta} $\ell=2$,
	\begin{align*}
		(\alpha_1,\beta_1,\kappa_1,\upsilon_1)=(0,\tilde{d},a,a+c),
		\quad \text{and} \quad
		({\alpha_2},{\beta_2},{\kappa_2},{\upsilon_2})=(0,\tilde{d},a,c),
	\end{align*}
	this condition is equivalent to the curve $\mathcal{X}_\vartheta$ having an affine $\F_{2^m}$-rational point $(x,z)$ with
	$x \notin \{ \tfrac{c}{a},\, 1+\tfrac{c}{a} \}$.
	Note that
	$\alpha_i\upsilon_i+\beta_i\kappa_i \neq 0$ for $i=1,2$
	since $ad \neq 0$.
	Moreover,
	$\upsilon_1/\kappa_1 \neq {\upsilon}_2/{\kappa}_2$ since $a\neq 0$.
	Therefore, the existence of such a point follows from
	Lemma~\ref{lem:curveDuv}.
	\hfill$\Box$\\[.5em]
	
	We now give a complete analysis of the generalized $\mathcal{PS}_{ap}$ bent function in the case $Q(x)=x$.
	That is, we consider the bent function
	\[
	f \colon \mathbb{F}_{2^{m}} \times \mathbb{F}_{2^{m}} \to \mathbb{F}_2,
	\qquad
	f(x,y)=\Tr_1^{m}\!\left(x\,y^{-\eta}\right).
	\]

	\begin{corollary}\label{cor:x}
		Let $m$, $k$, and $e$ be positive integers such that
		\[
		k \mid m, \qquad  e \equiv 2^{\ell} \mod{(2^{k}-1)},
		\qquad 
		\gcd(2^{m}-1,e)=1.
		\]
		Let $\eta$ denote the multiplicative inverse of $e$ modulo $2^{m}-1$.
		If $k \le \tfrac{m}{4}-3$, then the generalized $\mathcal{PS}_{ap}$ bent function $f(x,y)=\Tr_1^{m}\!\left(x\,y^{-\eta}\right)$ on $\mathbb{F}_{2^{m}} \times \mathbb{F}_{2^{m}}$
		satisfies the following properties:
		\begin{itemize}
			\item[(i)] The function $f$ admits a semibent decomposition on 
			$\mathcal{S}=\langle u,v\rangle^{\perp}$, for linearly independent
			$u,v \in \mathbb{F}_{2^{m}} \times \mathbb{F}_{2^{m}}$,
			if and only if 
			$u,v \in \{0\}\times \mathbb{F}_{2^{m}}$.
			
			\item[(ii)] The function $f$ does not admit any bent decomposition.
		\end{itemize}
	\end{corollary}

	\begin{proof} Let $u=(a,b)$ and $v=(c,d)$ be two linearly independent vectors in
		$\mathbb{F}_{2^{m}} \times \mathbb{F}_{2^{m}}$.
		By Proposition~\ref{mainprop}, the function $f$ does not admit any semibent or bent
		decomposition on $\mathcal{S}$ whenever $ad+bc \neq 0$.
		Therefore, we restrict our attention to the case $ad+bc = 0$.
		Recall that, since $u=(a,b)$ and $v=(c,d)$ are linearly independent vectors in
		$\mathbb{F}_{2^{m}} \times \mathbb{F}_{2^{m}}$, the condition $ad+bc = 0$
		is equivalent to $(c,d)=\lambda (a,b)$ for some $ \lambda \in \mathbb{F}_{2^{m}} \setminus \mathbb{F}_{2}$.
		
		By Lemma~\ref{lem:decom}, the function $f$ admits a semibent (respectively, bent)
		decomposition on $\mathcal{S}$ if and only if its dual $f^{\ast}(x,y)=\Tr_1^{m}\!\left(y^{2^\ell}\,x^{-2^\ell e}\right)$
		satisfies
		\[
		D_u D_v f^{\ast}=0
		\quad\text{(respectively, } D_u D_v f^{\ast}=1\text{).}
		\]
		Note that for the linear map $L(x,y)=(x^{2^\ell},y^{2^\ell})$ we have 
		$f^{\ast}(x,y)=\widetilde{f}(L(x,y))$, where 
		$\widetilde{f}(x,y)=\Tr_1^{m}(y\,x^{-e})$.
		Hence, by Lemma~\ref{lem:second}, it suffices to consider the second-order derivative of $\widetilde{f}$.
		
		By replacing $c=\lambda a$ and $d=\lambda b$,
		the second-order derivative $D_u D_v \widetilde{f}(x,y)$ can be expressed as
		\begin{align*}
			D_u D_v \widetilde{f}(x,y)
			=
			\Tr_1^{m}\!\Big(
			&(y+b(1+\lambda))\,(x+a(1+\lambda))^{-e}
			+(y+b)\,(x+a)^{-e} \\
			& \qquad \qquad \qquad \qquad\qquad +(y+\lambda b)\,(x+\lambda a)^{-e}
			+y\,x^{-e}
			\Big).
		\end{align*}
		
		We now proceed by a case-by-case analysis.
		
		\medskip
		
		\noindent \textbf{Case (i):}  $a=0$, and hence $b\neq 0$.  \\[.3em]
		In this case, we obtain $D_u D_v \widetilde{f}=0$ for all such choices of $u$ and $v$.
		Therefore, $ \widetilde{f}$ admits a semibent decomposition on $\langle u,v\rangle^{\perp}$.
		In particular, for any linearly independent 
		$u,v \in \{0\}\times \mathbb{F}_{2^{m}}$, 
		the function $\widetilde{f}$ admits a semibent decomposition on $\langle u,v\rangle^{\perp}$.
		Since 
		\[
		L^{-1}(\{0\}\times \mathbb{F}_{2^{m}})
		=
		\{0\}\times \mathbb{F}_{2^{m}},
		\]
		it follows from Lemma~\ref{lem:second} that 
		$f^{*}(x,y)$ admits a semibent decomposition on 
		$\langle u,v\rangle^{\perp}$ for all linearly independent 
		$u,v \in \{0\}\times \mathbb{F}_{2^{m}}$.
		
		\medskip
		
		\noindent \textbf{Case (ii):} $b=0$, and hence $a\neq 0$.  \\[.3em]
		Then applying the changes of variables $x\mapsto ax$ and $y\mapsto a^{e}y$ 
		and then setting $y=1$, we obtain
		\begin{align}\label{eq1}
			g(x)
			&=
			\Tr_1^{m}\!\Big(
			(x+1+\lambda)^{-e}
			+   (x+1)^{-e}
			+ (x+\lambda )^{-e}
			+ x^{-e}
			\Big).
		\end{align}
		Hence, if $D_u D_v f^{\ast}(x,y)$ is a constant function,
		then the function $g(x)$ in \eqref{eq1} must also be constant,
		which is impossible by Lemma~\ref{lem:curveDuv}. Then, by Lemma~\ref{lem:second}, the function $f^{*}$ does not admit neither a bent or a semibent decomposition on $\langle u,v\rangle^{\perp}$ for any linearly independent $u,v \in \mathbb{F}_{2^{m}} \times \{0\}$.
		
		\medskip
		
		\noindent \textbf{Case (iii):} $ab \neq 0$.  \\[.3em]
		Proceeding as before, we apply the change of variables $x \mapsto ax$ and
		$y \mapsto by$, and then set $\alpha = b a^{-e}$ and $y=0$. This yields
		\begin{align}\label{eq2}
			g(x)
			&=
			\Tr_1^{m}\!\Big(
			\alpha\Big(
			(1+\lambda)\,(x+1+\lambda)^{-e}
			+ (x+1)^{-e}
			+ \lambda (x+\lambda)^{-e}
			\Big)\,
			\Big).
		\end{align}
		As in the previous cases, if $D_u D_v \widetilde{f}(x,y)$ were a
		constant function, then $g(x)$ in~\eqref{eq2} would also
		have to be constant. This, however, is impossible by  Lemma~\ref{lem:curveDuv}. Then we similarly conclude that $f^{*}$ admits neither a bent nor a semibent decomposition on $\langle u,v\rangle^{\perp}$ for any such linearly independent elements $u,v$.
		
		\medskip
		
		This completes the proof.
	\end{proof}
	
	\begin{remark}
		Note that the generalized $\mathcal{PS}_{ap}$ bent function $f$ in
		Corollary~\ref{cor:x}, and hence its dual $f^{\ast}$, is also an $MM$ function.
		Consequently, $f$ and $f^{\ast}$ admit $m$-dimensional $\mathcal{M}$-subspaces,
		namely the canonical ones $\mathbb{F}_{2^{m}} \times \{0\}$ and
		$\{0\} \times \mathbb{F}_{2^{m}}$, respectively.
		It follows from Corollary~\ref{cor:x} that
		$\{0\} \times \mathbb{F}_{2^{m}}$ is the unique $m$-dimensional
		$\mathcal{M}$-subspace of $f^{\ast}$.
	\end{remark}
	
	We can similarly generalize Corollary~\ref{cor:x} to arbitrary permutation
	polynomials $Q(x)$ whose degree is sufficiently small compared to the
	cardinality of the underlying finite field. In particular, combined with
	Proposition~\ref{mainprop}, we obtain the following sufficient conditions under
	which the generalized $\mathcal{PS}_{ap}$ bent function admits neither a
	semibent nor a bent decomposition. The proof proceeds analogously to that of Corollary~\ref{cor:x}; therefore, we only give a sketch.
	
	\begin{theorem}\label{thm:gps}
		Let $m$, $k$, and $e$ be positive integers such that
		\[
		k \mid m, 
		\qquad  
		e \equiv 2^{\ell} \mod{(2^{k}-1)},
		\qquad 
		\gcd(2^{m}-1,e)=1.
		\]
		Let $\eta$ denote the multiplicative inverse of $e$ modulo $2^{m}-1$.
		Let $Q$ be a permutation polynomial of $\mathbb{F}_{2^{m}}$ 
		of odd degree at most $\tfrac{2^{m/4}-1}{4e}$ satisfying $Q(0)=0$.
		Define the generalized $\mathcal{PS}_{ap}$ bent function 
		$f \colon \mathbb{F}_{2^{m}} \times \mathbb{F}_{2^{m}} \to \mathbb{F}_2$ by $f(x,y)= \Tr_1^{m}\!\left(Q\!\left(x\,y^{-\eta}\right)\right)$.
		Then $f$ admits neither a semibent nor a bent decomposition provided that, 
		for every $\lambda \in \mathbb{F}_{2^{m}} \setminus \mathbb{F}_2$, 
		the Boolean function
		\begin{equation}\label{eq:Qgen}
			g_\lambda(x)
			= \Tr_1^{m}\!\left(
			Q\!\left((1+\lambda)x^{-e}\right)
			+ Q\!\left(x^{-e}\right)
			+ Q\!\left(\lambda x^{-e}\right)
			\right)
		\end{equation}
		on $\mathbb{F}_{2^{m}}$ is nonconstant.
	\end{theorem}

	\noindent{\it\textbf{Sketch of the proof.}}
	Similarly, by Proposition~\ref{mainprop}, it suffices to consider bent or semibent 
	decompositions on $\mathcal{S}=\langle u,v\rangle^{\perp}$,
	where $u=(a,b)$ and $v=(c,d)$ are linearly independent vectors in 
	$\mathbb{F}_{2^{m}} \times \mathbb{F}_{2^{m}}$ satisfying $ad+bc=0$.
	This condition implies that
	$(c,d) =\lambda (a,b)$ for some 
	$\lambda \in \mathbb{F}_{2^{m}}\setminus \mathbb{F}_{2}$.
	Moreover, by Lemma~\ref{lem:second}, it suffices to consider $\widetilde{f}(x,y)=\Tr_1^{m}(Q(y\,x^{-e}))$.
	
	\medskip
	
	\noindent
	\emph{Case $a=0$ (and hence $b\neq 0$).}
	Setting $y=0$ and applying the change of variables $x \mapsto \tilde{b}x$,
	where $\tilde{b}^{\,e}=b$, to $D_u D_v \widetilde{f}(x,y)$,
	we obtain the function
	$g_\lambda(x)$ in \eqref{eq:Qgen}. Hence, if $g_\lambda(x)$ is not constant, then
	$D_u D_v f^{\ast}(x,y)$ cannot be constant.
	
	\medskip
	
	\noindent
	\emph{Case $b=0$ (and hence $a\neq 0$).}
	Applying the changes of variables $x \mapsto ax$ and $y \mapsto a^{e}y$ to
	$D_u D_v \widetilde{f}(x,y)$ and then setting $y=1$, we obtain
	\[
	\Tr_1^{m}\!\Big(
	Q\!\left((x+1+\lambda)^{-e}\right)
	+ Q\!\left((x+1)^{-e}\right)
	+ Q\!\left((x+\lambda)^{-e}\right)
	+ Q\!\left(x^{-e}\right)
	\Big),
	\]
	which cannot be constant by Lemma~\ref{lem:curveDuv}. Consequently,
	$D_u D_v f^{\ast}(x,y)$ cannot be constant.
	
	\medskip
	
	\noindent
	\emph{Case $ab\neq 0$.}
	We apply the changes of variables $x \mapsto ax$ and $y \mapsto by$, and then set
	$\alpha = b a^{-e}$ and $y=0$. This yields
	\[
	\Tr_1^{m}\!\left(
	Q\!\left(\alpha (1+\lambda)(x+1+\lambda)^{-e}\right)
	+ Q\!\left(\alpha (x+1)^{-e}\right)
	+ Q\!\left(\alpha \lambda(x+\lambda)^{-e}\right)
	\right).
	\]
	Again, by Lemma~\ref{lem:curveDuv}, this expression cannot be constant, and hence
	$D_u D_v f^{\ast}(x,y)$ cannot be constant.
	
	\medskip
	
	The result now follows from Lemma~\ref{lem:decom}, which completes the proof.
	\hfill$\Box$
	
	\begin{example}
		Let $m$ be an odd integer and let $k$ be a positive integer with $\gcd(k,m)=1$. 
		Let $Q(x)=x^{2^k+1}$ be the Gold function. We recall that $Q$ is a permutation of $\mathbb{F}_{2^{m}}$ if and only if $m$ is odd and $\gcd(k,m)=1$. 
		
		For $\lambda \in \mathbb{F}_{2^{m}} \setminus \mathbb{F}_{2}$, the function $g_\lambda(x)$ defined in \eqref{eq:Qgen} satisfies
		\begin{align*}
			g_\lambda(x)
			&= \Tr_1^{m}\!\left(
			(1+\lambda)^{2^k+1} x^{-e(2^k+1)}
			+ x^{-e(2^k+1)}
			+ \lambda^{2^k+1} x^{-e(2^k+1)}
			\right) \\
			&= \Tr_1^{m}\!\left(
			\big((1+\lambda)^{2^k+1}+1+\lambda^{2^k+1}\big)
			x^{-e(2^k+1)}
			\right) \\
			&= \Tr_1^{m}\!\left(
			\lambda^{2^k}\, x^{-e(2^k+1)}
			\right).
		\end{align*}
		Since $\gcd(e(2^k+1),\,2^m-1)=1$, the function $g_\lambda$ is balanced for every nonzero $\lambda \in \mathbb{F}_{2^{m}}$. 
		
		Therefore, by Theorem~\ref{thm:gps}, the generalized $\mathcal{PS}_{ap}$ bent function 
		\[
		f \colon \mathbb{F}_{2^{m}} \times \mathbb{F}_{2^{m}} \to \mathbb{F}_2,
		\qquad 
		f(x,y)= \Tr_1^{m}\!\left(x^{2^k+1}\,y^{-\eta(2^k+1)}\right),
		\]
		where $\eta$ is the unique integer satisfying $\eta e \equiv 1 \mod{(2^{m}-1)}$, 
		does not admit any semibent or bent decomposition for all sufficiently small integers $k$ relative to $2^{m}$.
	\end{example}
	
	
	\section{Decomposition and concatenation of bent functions}
	\label{concat}
	
	Let $u, v$ be two linearly independent elements of $\V_n^{(2)}$, and let 
	$\mathcal{S} = \langle u, v \rangle^{\perp}$. 
	Since $\langle \cdot , \cdot \rangle_n$ is a nondegenerate inner product on 
	$\V_n^{(2)}$, we can write 
	$\V_n^{(2)} = \mathcal{S} \oplus \langle u, v \rangle$. 
	
	
	Let $f \colon \V_n^{(2)} \to \mathbb{F}_2$ be a Boolean function, and
	let $(f_1,f_2,f_3,f_4)$ be its decomposition with respect to
	$\mathcal{S}$. After a suitable change of coordinates as described
	above, we may identify $\mathcal{S}=\V_{n-2}^{(2)}$ and $u=(1,0)$, $v=(0,1)$.
	Thus, we may regard $f$ as a function on
	$\V_{n-2}^{(2)} \times \mathbb{F}_2 \times \mathbb{F}_2$ obtained by
	concatenation of $f_1, f_2, f_3, f_4 \colon \V_{n-2}^{(2)} \to \mathbb{F}_2$.
	That is, in this representation, the decomposition $(f_1,f_2,f_3,f_4)$
	satisfies
	\[
	f(x,y,z)=
	\begin{cases}
	f_1(x), & \text{if } (y,z)=(0,0),\\
	f_2(x), & \text{if } (y,z)=(0,1),\\
	f_3(x), & \text{if } (y,z)=(1,0),\\
	f_4(x), & \text{if } (y,z)=(1,1), 
	\end{cases}
	\]
	where $x\in\V_{n-2}^{(2)}$ and $y,z\in\mathbb{F}_2$.
	Equivalently,
	\[
	f(x,y,z)
	= f_1(x)
	+ yz \bigl(f_1 + f_2 + f_3 + f_4\bigr)(x)
	+ y \bigl(f_1 + f_3\bigr)(x)
	+ z \bigl(f_1 + f_2\bigr)(x).
	\]
	In other words, $f$ is obtained by concatenation of the Boolean functions
	$f_1, f_2, f_3, f_4 $.
	
	We define the $4$-concatenation (or simply, concatenation) of arbitrary Boolean functions
	$f_1, f_2, f_3, f_4 \colon \V_n^{(2)} \to \mathbb{F}_2$ as the Boolean
	function $f = f_1 \,\|\, f_2 \,\|\, f_3 \,\|\, f_4$ from
	$\V_n^{(2)} \times \mathbb{F}_2^2$ to $\mathbb{F}_2$ given by
	\[
	f(x,y,z)
	= f_1(x)
	+ yz \bigl(f_1 + f_2 + f_3 + f_4\bigr)(x)
	+ y \bigl(f_1 + f_3\bigr)(x)
	+ z \bigl(f_1 + f_2\bigr)(x),
	\]
	where $x \in \V_n^{(2)}$ and $y,z \in \mathbb{F}_2$. 
	
	We remark that the concatenation method has been used efficiently for
	the secondary construction of Boolean bent functions outside the \(MM^{\#}\) class. For instance, see
	\cite{Kudin25,PaPoKuZ,PasalicEtAl2023} as well as the survey paper
	\cite{Pasalic26} and the references therein.
	The necessary and sufficient condition for the concatenation of bent
	functions to be bent is given as follows.

	\begin{lemma}[{\cite[Theorem~III.1]{hpz}}]
		\label{ffff=1}
		Let $f_1,f_2,f_3,f_4$ be four Boolean bent functions. Then the concatenation $f = f_1 \,\|\, f_2 \,\|\, f_3 \,\|\, f_4$
		is bent if and only if $f_1^* + f_2^* + f_3^* + f_4^* = 1$.
	\end{lemma}
	
	We recall that a $GMM$ function $f \colon \mathbb{F}_{2^n} \times \mathbb{F}_{2^k}^2 \to \mathbb{F}_2$
	is defined by $f(x,y,z) = f^{(z)}(x) + \Tr_1^k(yz)$,
	where, for each $z \in \mathbb{F}_{2^k}$, the function
	$f^{(z)} \colon \mathbb{F}_{2^n} \to \mathbb{F}_2$ is bent.
	In the special case $k=1$, that is, for
	$f \colon \mathbb{F}_{2^n} \times \mathbb{F}_2 \times  \mathbb{F}_2 \to \mathbb{F}_2$
	given by $f(x,y,z)=f^{(z)}(x)+yz$, we observe that
	$f(x,y,z)=f^{(0)}(x)$ whenever $z=0$, whereas for $z=1$ we have
	$f(x,y,z)=f^{(1)}(x)$ if $y=0$ and
	$f(x,y,z)=f^{(1)}(x)+1$ if $y=1$.
	Consequently, as noted in \cite{akkmpp26}, $f$ can be expressed as the
	concatenation $f = f^{(0)} \,\|\, f^{(1)} \,\|\, f^{(0)} \,\|\, \bigl(f^{(1)}+1\bigr)$. 
	In this case, the necessary and sufficient condition on the dual
	functions in Lemma~\ref{ffff=1} is trivially satisfied.
	More generally, as in the case $k=1$, $GMM$ functions may be viewed as
	concatenations of bent functions over $\mathbb{F}_{2^n}$.
	
	In general, concatenations of bent functions (particularly the ones of the form $f = f^{(0)} \,\|\, f^{(1)} \,\|\, f^{(0)} \,\|\, \bigl(f^{(1)}+1\bigr)$) have been studied for constructing bent functions outside the
	completed $MM$ class. In the recent paper \cite{akkmpp26}, the
	concatenation of vectorial $\mathcal{PS}_{ap}$ bent functions is
	investigated. The motivation stems from the fact that the components of
	a vectorial $\mathcal{PS}_{ap}$ bent function $F \colon \V_n^{(2)} \to \V_m^{(2)}$
	satisfy the following property: for any nonzero
	$\alpha, \beta \in \V_m^{(2)}$ with $\alpha \ne \beta$, we have
	\begin{align}\label{eq:comp}
		(F_\alpha)^{*} + (F_\beta)^{*} = (F_{\alpha+\beta})^{*}.  
	\end{align}
	
	The property in \eqref{eq:comp} holds for almost all bent functions arising from bent partitions; in particular, it applies to vectorial generalized $\mathcal{PS}_{ap}$ bent functions $F \colon \V_n^{(2)} \to \V_m^{(2)}$. Owing to their explicit representation, analogous to that of $\mathcal{PS}_{ap}$ bent functions, one can construct bent functions by concatenating the components of $F$ as $f = F_\alpha \,\|\, F_\beta \,\|\, F_\gamma \,\|\, F_{\alpha+\beta+\gamma} + 1$, where $\alpha, \beta, \gamma \in \V_m^{(2)}$ are nonzero elements such that $\alpha + \beta + \gamma \ne 0$, as stated below. The proof proceeds by analogous arguments and is therefore omitted.
	
	\begin{theorem}
		Let $m$, $k$ be integers such that $k$ divides $m$ and $\gcd(2^m-1, 2^k+1) = 1$. Set
		$e =2^k+1$.
		Let $P$ be any permutation of $\F_{2^k}$, 
		and let $\alpha,\beta,\gamma$ be any nonzero elements of $\F_{2^k}$ such that $\alpha + \beta + \gamma \ne 0$. Then 
		\begin{equation}
			\label{PSffff}
			f(x,y,z_1,z_2)
			= \Tr_1^k \!\left(
			\bigl((1+z_1+z_2)\alpha + z_2\beta + z_1\gamma\bigr)
			\, P\!\left(\Tr_k^m\!\left(yx^{-e}\right)\right)
			\right)
			+ z_1 z_2.
		\end{equation}
		is a bent function from $\F_{2^m}\times\F_{2^m}\times\F_2\times\F_2$ to $\F_2$.
	\end{theorem}
	
	We conclude this section with a more general construction of bent
	functions arising from the structure of the preimage distribution of
	$f$ given in \eqref{PSffff}. Without loss of generality, we may assume
	that $P$ is a permutation of $\F_{2^m}$ satisfying $P(0)=0$. Under this
	assumption, the preimage distribution of $f$ in \eqref{PSffff} is given below.
	
	Note that the collection of subsets
	\[
	\bigl\{
	U \times \{(z_1,z_2)\}
	\bigr\}
	\;\cup\;
	\bigl\{
	\mathcal A(\gamma) \times \{(z_1,z_2)\}
	:\, \gamma \in \F_{2^k}
	\bigr\}
	\]
	forms a disjoint cover of $\F_{2^m} \times \F_{2^m} \times \F_2 \times \F_2$,
	where $U$ and $\mathcal A(\gamma)$ are defined in \eqref{eq:U} and \eqref{eq:Agamma}, respectively.
	The subsets $U\times\{(0,0)\}$, $U\times\{(0,1)\}$, and 
	$U\times\{(1,0)\}$ are mapped to $0$, whereas 
	$U\times\{(1,1)\}$ is mapped to $1$. 
	Moreover, exactly half of the sets 
	$\{\mathcal{A}(\gamma)\times\{(y,z)\} \,:\, \gamma \in \F_{2^k}\}$ 
	are mapped to $0$; consequently, the remaining $2^{k-1}$ are mapped to $1$ whenever $(y,z)\in\{(0,0),(1,0),(0,1)\}$. 
	Finally, exactly half of the sets 
	$\{\mathcal{A}(\gamma)\times\{(1,1)\} \,:\, \gamma \in \F_{2^k}\}$ 
	are mapped to $0$ in such a way that, for each $\gamma \in \F_{2^k}$, if 
	$\mathcal{A}(\gamma)\times\{(0,0)\}$, 
	$\mathcal{A}(\gamma)\times\{(0,1)\}$, and 
	$\mathcal{A}(\gamma)\times\{(1,0)\}$ 
	are mapped to $a_0$, $a_1$, and $a_2$, respectively, then 
	$\mathcal{A}(\gamma)\times\{(1,1)\}$ is mapped to 
	$a_0 + a_1 + a_2 + 1$.
	
	Motivated by the preimage distribution of $f$ in \eqref{PSffff}, we present the following construction of bent functions, analogous to that obtained from the Desarguesian spread in \cite{akkmpp26}. Although the arguments are more involved, they are similar to those in \cite{akkmpp26}; hence, we omit the proof.
	
	\begin{theorem}
		\label{thmparti}
		Let $\Omega_1=\{U,\mathcal{A}(\gamma)\,:\,\gamma\in\F_{2^k}\}$ 
		be the generalized Desarguesian spread of 
		$\F_{2^m}\times\F_{2^m}$ given in \eqref{eq:Gamma}. 
		This induces a partition of 
		$\F_{2^m}\times\F_{2^m}\times\F_2\times\F_2$ 
		into the sets $U\times\{(y,z)\}$ and 
		$\mathcal{A}(\gamma)\times\{(y,z)\}$, 
		where $\gamma\in\F_{2^k}$ and $(y,z)\in\F_2^2$. 
		
		Assign to each quadruple $(a_0,a_1,a_2,a_3)\in\F_2^4$ 
		with an odd number of $1$'s exactly $2^{k-3}$ sets 
		$\mathcal{A}(\gamma)$. Define 
		$f:\F_{2^m}\times\F_{2^m}\times\F_2\times\F_2\to\F_2$ 
		to be constant on each part of this partition as follows:\\
		$f=0$ on $U\times\{(y,z)\}$ for $(y,z)\in\{(0,0),(0,1),(1,0)\}$, 
		and $f=1$ on $U\times\{(1,1)\}$. 
		If $\mathcal{A}(\gamma)$ is assigned to $(a_0,a_1,a_2,a_3)$, 
		then $f(x,y,z_1,z_2)=a_i$ whenever 
		$(x,y)\in\mathcal{A}(\gamma)$ and 
		$(z_1,z_2)=(0,0),(0,1),(1,0),(1,1)$ 
		for $i=0,1,2,3$, respectively. 
		Then $f$ is a bent function.
	\end{theorem}

	
	\section*{Acknowledgements}
	
	The initial work on this project began during the “Women in Numbers Europe~5 (WINE-5)” workshop, held at the University of Split in August 2025. The authors are grateful to the University of Split and the supporting institutions for making this conference and the resulting collaboration possible. They would especially like to thank the organizers of WINE-5 —Marcela Hanzer, Borka Jadrijević, Pınar Kılıçer, and Lejla Smajlović—for their dedication and hard work, which made the meeting both exceptionally fruitful and enjoyable.
	
	This study was supported by the Scientific and Technological Research  Council of Turkey (TÜBİTAK) under Grant Number 125F396. S.~A. and N.~A. thank TÜBİTAK for its support.
	T.~K. is supported by the FWF Project~P~35138.
	B.~T. is supported by the Swiss National Foundation through grant no.~212865.
	
	\begin{thebibliography}{99}
		
		
		\bibitem{aakm25} S. Alkan, N. Anbar, T. Kalayc\i, W. Meidl, Bent partition, vectorial dual-bent function, and LP-Packing constructions. IEEE Trans. Inform. Theory 71 (2025), 752--767. 
		
		\bibitem{akkmpp26} N.~Anbar, T.~Kalayc\i, S.~Kudin, W.~Meidl, E.~Pasalic, A.~Polujan,
		{Concatenation and decomposition of \(\mathcal{PS}\) bent functions and related \(\mathcal{M}\)-subspaces}. Preprint 2026. 
		
		\bibitem{akm25} N. Anbar, T. Kalayc\i, W. Meidl, A survey on bent partitions. To appear in Cryptogr. Commun.
		
		\bibitem{Anbar2025Preprint}
		N.~Anbar, F.-W.~Fu, T.~Kalaycı, W.~Meidl, J.~Wang, Y.~Wei,
		Analysis of some classes of bent partitions and vectorial bent functions. To appear in Des. Codes Cryptogr.
		
		\bibitem{akm24} N.~Anbar, T.~Kalayc{\i}, W.~Meidl,
		Amorphic association schemes from bent partitions.
		Discrete Math. 347 (2024), Paper No.~113658, 13~pp.
		
		\bibitem{akm23} N.~Anbar, T.~Kalaycı, and W.~Meidl,
		On generalized spread bent partitions.
		Cryptogr. Commun. 15 (2023), 1217--1234.
		
		
		\bibitem{akm22a} N. Anbar, T. Kalayc\i, W. Meidl, Generalized semifield spreads. Des. Codes Cryptogr. 91 (2023), 545--562. 
		
		\bibitem{nuwi22} N. Anbar, W. Meidl, Bent partitions. Des. Codes Cryptogr. 90 (2022), 1081--1101. 
		
		\bibitem{Camion92} P.~Camion, C.~Carlet, P.~Charpin, N.~Sendrier, On correlation-immune functions.
		In Advances in Cryptology --- CRYPTO~'91, J.~Feigenbaum, Ed., 86--100.
		Springer, Berlin, Heidelberg, 1992.
		
		\bibitem{CaCha} A. Canteaut, P. Charpin, Decomposing bent functions. IEEE Trans. Inform. Theory 49
		(2003), 2004--2019.
		
		\bibitem{Carlet2021} C.~Carlet, Boolean Functions for Cryptography and Coding Theory. Cambridge University Press, Cambridge, 2021.
		
		\bibitem{Carlet2004} C.~Carlet,
		On the confusion and diffusion properties of Maiorana--McFarland’s
		and extended Maiorana--McFarland’s functions.
		J.\ Complexity 20 (2004), 182--204 (2004).
		
		\bibitem{cmpGMM} A. \c Ce\c smelio\u glu, W. Meidl, A. Pott, Generalized Maiorana--McFarland class and normality of $p$-ary bent functions.
		Finite Fields Appl. 24 (2013), 105--117.
		
		\bibitem{dillon} J.~F.~Dillon, Elementary Hadamard difference sets,
		Ph.D. dissertation, University of Maryland, 1974.
		
		\bibitem{gold} R. Gold, Maximal recursive sequences with $3$-valued recursive cross--correlation functions. IEEE Trans. Inform. Theory 14 (1968), 154--156.
		
		\bibitem{hkt} J.~W.~P. Hirschfeld, G.~Korchm{\'a}ros, F.~Torres,
		Algebraic Curves over a Finite Field.
		Princeton University Press, 2013.
		
		\bibitem{hpz} S. Hodzic, E. Pasalic, W. Zhang, Generic constructions of five-valued spectra 
		boolean functions. IEEE Trans. Inform. Theory 65 (2019), 7554--7565.
		
		\bibitem{Kudin25}
		S.~Kudin, E.~Pasalic, A.~Polujan, F.~Zhang,
		The algebraic characterization of $\mathcal{M}$-subspaces of bent concatenations and its application. IEEE Trans. Inform. Theory 71 (2025), 3999--4011.

		\bibitem{KSW85} P.~V.~Kumar, R.~A.~Scholtz, L.~R.~Welch, Generalized bent functions and their properties. J. Combin.\ Theory Ser.~A {40} (1985), 90--107.
		
		\bibitem{ll} P.~Langev, G.~Leander, Counting all bent functions in dimension eight,
		Des. Codes Cryptogr. 59 (2011), 193--205.
		
		\bibitem{Letal2025}
		Y.~Li, H.~Kan, S.~Mesnager, J.~Peng, L.~Zheng,
		Direct approaches for generic constructions of plateaued functions
		and bent functions outside $\mathcal{M}^\#$.
		IEEE Trans. Inform. Theory 71 (2025), 1400--1418.
		
		\bibitem{mm} R.~L.~McFarland, A family of difference sets in non-cyclic groups.
		J. Combin. Theory Ser.~A 15 (1973), 1--10.
		
		\bibitem{mp} W.~Meidl, I.~Pirsic, Bent and $\mathbb{Z}_{2^k}$-bent functions from spread-like partitions.
		Des.\ Codes Cryptogr. 89 (2021), 75--89.
		
		\bibitem{Pasalic26}
		E.~Pasalic, A.~Polujan, F.~Zhang, S.~Kudin, Bent functions and the completed Maiorana--McFarland class.
		Preprint, 2026.
		
		\bibitem{PaPoKuZ} E. Pasalic, A. Polujan, S. Kudin, F. Zhang, Design and Analysis of bent functions
		using $\mathcal{M}$-subspaces. IEEE Trans. Inform. Theory 70 (2024), 4464--4477.
		
		\bibitem{PasalicEtAl2024}
		E.~Pasalic, A.~Bapic, F.~Zhang, Y.~Wei,
		Using $P_{\tau}$ property for designing bent functions provably
		outside the completed Maiorana--McFarland class.
		Des.\ Codes Cryptogr. 92 (2024), 2639--2654.
		
		\bibitem{PasalicEtAl2023}
		E.~Pasalic, A.~Bapic, F.~Zhang, Y.~Wei,
		Explicit infinite families of bent functions outside the completed
		Maiorana--McFarland class.
		Des.\ Codes Cryptogr. 91 (2023), 2365--2393.
		
		\bibitem{pp} A.~Polujan, A.~Pott, Cubic bent functions outside the completed Maiorana--McFarland class. Des. Codes Cryptogr. 88 (2020), 1701--1722.
		
		\bibitem{sti}
		H.~Stichtenoth, Algebraic Function Fields and Codes.
		2nd~ed.,
		Graduate Texts in Mathematics, vol.~254,
		Springer-Verlag, 2009.
		
		\bibitem{wfw} J.~Wang, F.-W.~Fu, Y.~Wei, Bent partitions, vectorial dual-bent functions and partial difference sets. IEEE Trans. Inform. Theory 69 (2023), 7414--7425.
		
		\bibitem{wfq}  G.~Weng, R.~Feng, W.~Qiu, On the ranks of bent functions.
		Finite Fields Appl. 13 (2007), 1096--1116.
	
	\end{thebibliography}
	
	
	
\end{document}

